\documentclass{article}
\usepackage{readmasters}
\begin{document}
\origpage{403}
\section*{Ueber algebraische Gebilde mit eindeutigen Transformationen in sich.}

Von \emph{A. Hurwitz} in Königsberg i.~Pr.

Die vorliegende Abhandlung habe ich in drei Abschnitte eingetheilt, über deren Inhalt ich hier einen kurzen Ueberblick voraufschicke.

Der erste Abschnitt bringt zunächst einen neuen Beweis des Satzes, nach welchem eine Riemann'sche Fläche (ein irreducibles algebraisches Gebilde) nur eine endliche Zahl von eindeutigen Transformationen in sich besitzen kann, wenn das Geschlecht $p$ der Fläche grösser als $1$ ist. Einen wesentlichen Stützpunkt dieses Beweises bildet die nähere Untersuchung derjenigen besonderen Stellen auf einer Fläche vom Geschlechte $p$, in welchen eine eindeutige algebraische Function von niedrigerer als der $(p+1)^{\text{ten}}$ Ordnung unendlich wird, ohne noch an anderen Stellen unstetig zu werden${}^{*)}$.

Die Anzahl dieser besonderen Stellen beträgt stets $p(p^2 - 1)$. Jedoch ist dabei vorausgesetzt, dass jede Stelle mit einer gewissen Multiplicität gezählt wird. Ich zeige nun, dass die Anzahl jener Stellen, wenn man jede Stelle einfach rechnet, also das Wort „Anzahl“ in seiner ursprünglichen Bedeutung versteht, stets grösser ist als $2p+2$. Eine Ausnahme bildet der hyperelliptische Fall, in welchem die Anzahl gleich $2p+2$ ist.

Schliesst man den hyperelliptischen Fall aus, so sind die erwähnten Stellen auf der Curve der $\varphi$ im Raume von $p-1$ Dimensionen mit denjenigen Punkten identisch, in welchen eine Ebene hyperosculirt, d.~i. mindestens $p$-punktig berührt. Ich bespreche kurz die allgemeinere Frage nach denjenigen Punkten auf der Curve der $\varphi$, in welchen eine

\textbf{*)} Diese Stellen spielen auch in den älteren Beweisen des Satzes eine Rolle. Eine Zusammenstellung der sich auf den Satz beziehenden Litteratur habe ich in Nr.~3 meiner Note: „Ueber diejenigen algebraischen Gebilde, welche eindeutige Transformationen in sich zulassen“, gegeben. Mathematische Annalen, Bd.~32 oder Göttinger Nachrichten aus dem Jahre 1887.

\origpage{404}
Fläche $k^{\text{ter}}$ Ordnung, wo $k > 1$ ist, hyperosculirt, d.~i. mindestens $(2k-1)(p-1)$-punktig berührt. Die Anzahl dieser Stellen ist stets gleich
\[
(2k-1)^2 \cdot p \cdot (p-1)^2,
\]
wenn man jede Stelle mit einer gewissen Multiplicität rechnet, welche ich genau definire.

Der zweite Abschnitt giebt die Construction der Riemann'schen Flächen, welche eine endliche Gruppe von eindeutigen Transformationen in sich besitzen, aus unabhängigen Bestimmungsstücken. Man erhält die allgemeinste derartige Fläche, wenn man eine beliebig gewählte Fläche $\Phi$ in $r$ Exemplaren nimmt, diese Exemplare auf einander legt und in geeigneter Weise zu einer einzigen Fläche mit einander verbindet. Die so entstehenden Flächen sind keine anderen, als die von Herrn \emph{Klein} eingeführten und von Herrn \emph{Dyck} nach verschiedenen Richtungen untersuchten „regulären“ Riemann'schen Flächen, welche die Verzweigung von Galois'schen Resolventen darstellen${}^{*)}$.

Ist eine Riemann'sche Fläche $F$ vom Geschlecht $p$ $r$-blättrig über der Fläche $\Phi$ vom Geschlecht $\pi$ ausgebreitet, so besteht die Gleichung
\[
2p - 2 = W + r(2\pi - 2),
\]
wo $W$ die Gesammtzahl der Verzweigungen von $F$ in Bezug auf $\Phi$ bedeutet.

Aus dieser Gleichung leite ich beiläufig die von Herrn \emph{Zeuthen} aufgestellte Relation zwischen den Geschlechtszahlen zweier mehrdeutig aufeinander bezogener algebraischen Gebilde ab. Eben diese Gleichung führt zu einer wesentlichen Ergänzung des im ersten Abschnitte bewiesenen Satzes. Es ergiebt sich nämlich, dass eine Riemann'sche Fläche vom Geschlecht $p > 1$ nicht mehr als $84(p-1)$ eindeutige Transformationen in sich besitzen kann.

Die Untersuchungen des dritten Abschnitts beziehen sich auf die Integrale erster Gattung einer Riemann'schen Fläche ($p > 1$), von welcher vorausgesetzt wird, dass sie eine Gruppe eindeutiger Transformationen in sich besitzt.

Jeder eindeutigen Transformation der Fläche in sich entspricht eine lineare Transformation der Integrale erster Gattung, welche man auch in die Gestalt einer homogenen linearen Transformation der Differentiale erster Gattung setzen kann. Der Gruppe eindeutiger Transformationen der Fläche in sich entspricht eine Gruppe von homogenen linearen Transformationen der Differentiale erster Gattung, und ich beweise zuerst, dass beide Gruppen holoedrisch isomorph auf einander bezogen sind.

\textbf{*)} Vergleiche die unten folgenden Citate.

\origpage{405}
Es handelt sich dann weiter darum, die Wurzeln der charakteristischen Gleichung für die einzelne homogene lineare Transformation der Differentiale erster Gattung zu bestimmen. Die Bestimmung geschieht auf Grund der Theorie gewisser Functionensysteme auf einer Riemann'schen Fläche. Die Functionen eines solchen Systems sind wesentlich durch die Eigenschaft charakterisirt, dass sie sich auf geschlossenen Wegen bis auf multiplicative Constanten reproduciren. Die eindeutigen algebraischen Functionen bilden einen speciellen Fall eines solchen Functionensystems. Sie entsprechen dem Falle, in welchem jene multiplicativen Constanten sämmtlich gleich $1$ sind. Auf die Theorie der allgemeinen Functionensysteme, welche an sich von Interesse erscheint, hoffe ich demnächst zurückzukommen.

Ich bemerke noch, dass die Resultate des dritten Abschnittes mannigfache Anwendungen auf die von Herrn \emph{Klein} entworfene Theorie der elliptischen Modulfunctionen gestatten. Thatsächlich bin ich auch durch Untersuchungen über Classenzahlrelationen${}^{*)}$, welche sich auf jene Theorie stützen, zu den hier behandelten allgemeineren Fragen hingeführt worden.

Schliesslich will ich noch in Betreff des allgemeinen Begriffes einer Riemann'schen Fläche auf die folgenden Publicationen verweisen: \emph{F.~Klein}, „Ueber Riemann's Theorie der algebraischen Functionen und ihrer Integrale“ (Leipzig 1882) und „Neue Beiträge zur Riemann'schen Functionentheorie“ Mathematische Annalen, Bd.~21. \emph{R.~Dedekind} und \emph{H.~Weber}: „Theorie der algebraischen Functionen einer Veränderlichen“. Crelle's Journal Bd.~92.

Soweit es sich nur um Lagebeziehungen, also um Wege und Schnitte auf einer Fläche handelt, denke ich mir dieselbe stets als eine frei im Raume liegende geschlossene Ringfläche mit $p$ Oeffnungen.

\section*{I. Abschnitt.}

\subsection*{1.}

\emph{Eine Riemann'sche Fläche vom Geschlecht $p$ besitze eine eindeutige Transformation $S$ in sich. Dann beträgt die Anzahl der verschiedenen Stellen $P$ der Fläche, welche mit ihrer entsprechenden Stelle $P'$ zusammenfallen, höchstens $2p+2$, es sei denn, dass $S$ die identische Transformation ist, d.~h. dass jede Stelle $P$ mit ihrer entsprechenden $P'$ zusammenfällt.}

In der That, angenommen $S$ sei nicht die identische Transformation und $Q$ irgend eine Stelle, welche nicht mit ihrer entsprechenden $Q'$

\textbf{*)} „Ueber die Classenzahlrelationen und Modularcorrespondenzen primzahliger Stufe“ Berichte der k.~sächsischen Gesellschaft der Wissenschaften vom 4.~Mai 1885.

\origpage{406}
zusammenfällt. Wir betrachten dann eine eindeutige algebraische Function des Ortes auf unserer Fläche, welche nur bei $Q$ und zwar von einer Ordnung $\mu \leqq p+1$ unendlich wird. Eine solche Function existirt bekanntlich immer. Sind nun $z$ und $z'$ die Werthe, welche diese Function in den entsprechenden Stellen $P$ und $P'$ annimmt, so können wir die Differenz $z' - z$ als eine Function der Stelle $P$ ansehen. Diese Function ist eine eindeutige algebraische Function der Fläche, welche $2\mu$-Mal unendlich wird, nämlich $\mu$-fach an der Stelle $Q$ und $\mu$-fach an der Stelle $Q_1$, deren entsprechende die Stelle $Q$ ist. Folglich wird $z' - z$ auch $2\mu$-Mal Null. Nun ist aber jede Stelle $P$, welche mit ihrer entsprechenden Stelle $P'$ coincidirt, eine Nullstelle von $z' - z$. Folglich ist die Anzahl dieser Stellen höchstens $2\mu \leqq 2p+2$, w.~z.~b.~w.

\emph{Wenn die Riemann'sche Fläche hyperelliptisch ist, so beträgt die im vorigen Satze erwähnte Anzahl höchstens $4$, es sei denn, dass die Transformation $S$ die identische ist, oder diejenige, welche je zwei Stellen einander zuordnet, in denen die zweiwerthige Function der Fläche denselben Werth annimmt.}

Sind nämlich $z$ und $z'$ die Werthe dieser Function in zwei entsprechenden Stellen $P$ und $P'$, so ist $z' - z$ aufgefasst als Function der Stelle $P$ eine eindeutige algebraische Function, welche höchstens $4$~Mal unendlich wird. Folglich wird $z' - z$ auch höchstens $4$~Mal Null, es sei denn, dass $z' - z$ identisch Null ist. Hieraus folgt ohne Weiteres die Richtigkeit unseres Satzes.

\emph{Eine Riemann'sche Fläche, deren Geschlecht $p$ grösser als $1$ ist, kann nur eine endliche Anzahl von eindeutigen Transformationen in sich besitzen.}

Zur Abkürzung werde ich mich der Ausdrucksweise bedienen „die Transformation $S$ ersetze den Punkt $P$ durch den Punkt $P'$“, wenn dem Punkte $P$ vermöge $S$ der Punkt $P'$ entspricht. Ferner werde ich, wie üblich, unter $S^{-1}$ diejenige Transformation verstehen, welche den Punkt $P'$ durch den Punkt $P$ ersetzt. Endlich möge, wenn $S$ den Punkt $P$ durch $P'$ und $T$ den Punkt $P'$ durch $P''$ ersetzt, unter $ST$ diejenige Transformation verstanden werden, welche den Punkt $P$ durch $P''$ ersetzt.

Dies vorausgeschickt, wende ich mich zu dem Beweise des aufgestellten Satzes und betrachte zuerst den Fall, wo die Riemann'sche Fläche nicht hyperelliptisch ist.

In diesem Falle giebt es, wie ich in der folgenden Nummer zeigen werde, stets eine Gruppe von $r > 2p+2$ Stellen $P_1, P_2, \dots P_r$, welche sich bei jeder eindeutigen Transformation der Fläche in sich nur untereinander vertauschen können. Hieraus schliessen wir sofort, dass die Fläche eine endliche Anzahl, nämlich höchstens $r!$ Transformationen in sich besitzen kann. Denn zwei Transformationen

\origpage{407}
$S$ und $S'$, welche dieselbe Vertauschung der Stellen $P_1, P_2, \dots P_r$ hervorbringen, sind nothwendig identisch, da $S'S^{-1}$ die $r$ Stellen fest lässt und also, nach dem oben Bewiesenen, die identische Transformation ist.

Wenn nun zweitens die Riemann'sche Fläche hyperelliptisch ist, so sei $z$ die zweiwerthige Function der Fläche. Danu\ednote{Printed Danu instead of Dann; an apparent compositor misprint with inverted n.} müssen sich die $2p+2$ Stellen, an welchen $dz$ von der zweiten Ordnung verschwindet, bei jeder eindeutigen Transformation der Fläche in sich unter einander vertauschen. Sind daher $S$ und $S'$ zwei solche Transformationen, welche dieselbe Vertauschung dieser Stellen hervorbringen, so bleiben bei der Transformation $S'S^{-1}$ mehr als $4$, nämlich die erwähnten $2p+2$ Stellen fest. Folglich ist $S'S^{-1}$ entweder die identische Transformation oder diejenige, welche je zwei Stellen einander zuordnet, in denen die zweiwerthige Function denselben Werth annimmt. Hieraus folgt nun sofort, dass die Fläche höchstens doppelt so viele Transformationen in sich besitzen kann, als die $2p+2$ Stellen unter einander vertauschbar sind, also höchstens $2(2p+2)!$.

\subsection*{2.}

Es sei $F$ eine Riemann'sche Fläche vom Geschlecht $p > 1$, ferner seien $u_1, u_2, \dots u_p$ zugehörige linear unabhängige Integrale 1.~Gattung. Bezeichnet nun $u$ eine beliebige Function des Ortes auf der Fläche $F$, so besitzt die Determinante
\[
\Delta_u =
\begin{vmatrix}
\frac{du_1}{du}, & \frac{du_2}{du}, & \dots & \frac{du_p}{du} \\
\frac{d^2u_1}{du^2}, & \frac{d^2u_2}{du^2}, & \dots & \frac{d^2u_p}{du^2} \\
\cdot & \cdot & \dots & \cdot \\
\cdot & \cdot & \dots & \cdot \\
\frac{d^p u_1}{du^p}, & \frac{d^p u_2}{du^p}, & \dots & \frac{d^p u_p}{du^p}
\end{vmatrix}
\tag{1}
\]
folgende Eigenschaften. Sie verschwindet nicht identisch, da $u_1, u_2, \dots u_p$ linear unabhängig sind. Sie multiplicirt sich mit einer nicht verschwindenden Constanten, wenn $u_1, u_2, \dots u_p$ durch irgend ein anderes System von $p$ überall endlichen Integralen ersetzt werden. Endlich ist
\[
\Delta_u = \Delta_t \cdot \left(\frac{dt}{du}\right)^{\frac{p(p+1)}{2}},
\tag{2}
\]
wenn $t$ irgend eine andere Function des Ortes auf der Fläche $F$ bezeichnet.

\origpage{408}
Wir wollen nun für $u$ insbesondere ein überall endliches Integral wählen. Dann stellt die Determinante $\Delta_u$ eine eindeutige algebraische Function der Fläche $F$ vor. Um diese näher zu untersuchen, betrachten wir eine Stelle $P$ der Fläche und bezeichnen mit $t$ eine Function des Ortes, welche in $P$ von der ersten Ordnung verschwindet.

Die Gleichung (2) zeigt nun, dass $\Delta_u$ nur an den $2p - 2$ Nullstellen von $du$ unendlich wird und zwar an jeder von der Ordnung $\frac{p(p+1)}{2}$. Die Gesammtordnung des Unendlichwerdens beträgt also $(p-1)p(p+1)$ und ebenso gross ist also auch die Gesammtordnung des Verschwindens. Nun wird aber $\frac{dt}{du}$ an keiner Stelle Null; also folgt:

„\emph{Es giebt stets eine endliche Zahl von Stellen $P$, an welchen die Determinante $\Delta_t$ verschwindet. Sind $P_1, P_2, \dots P_r$ diese Stellen und $m_1, m_2, \dots m_r$ die zugehörigen Ordnungszahlen des Verschwindens von $\Delta_t$, so ist}
\[
m_1 + m_2 + \dots + m_r = (p-1)p(p+1).
\tag{3}
\]
Es handelt sich jetzt um eine nähere Untersuchung der Zahlen $m_1, m_2, \dots m_r$.

Zu dem Ende betrachten wir eine beliebige Stelle $P$ der Fläche und wählen die überall endlichen Integrale so, dass ihre Entwicklungen an der Stelle $P$ die Gestalt haben
\[
\begin{aligned}
u_1 &= t^{\varrho_1} + \dots, \\
u_2 &= t^{\varrho_2} + \dots, \\
\cdot & \cdot \dots \cdot \\
u_p &= t^{\varrho_p} + \dots,
\end{aligned}
\]
wobei die Exponenten der Anfangsglieder den Bedingungen
\[
0 < \varrho_1 < \varrho_2 < \dots < \varrho_p
\tag{4}
\]
genügen. Diese Wahl der Integrale ist, wie leicht zu sehen, stets und nur auf eine Weise möglich. Die Entwicklung von $\Delta_t$ beginnt dann an der Stelle $P$ mit dem Gliede
\[
c \cdot t^{\varrho_1 + \varrho_2 + \dots + \varrho_p - \frac{p(p+1)}{2}} = c \cdot t^m,
\]
wenn wir zur Abkürzung
\[
m = \varrho_1 + \varrho_2 + \dots + \varrho_p - \frac{p(p+1)}{2}
\tag{5}
\]
setzen. Die Zahlen $\varrho_1, \varrho_2, \dots, \varrho_p$ besitzen eine einfache Bedeutung. Man betrachte nämlich die eindeutigen algebraischen Functionen der Fläche, welche nur an der Stelle $P$ unendlich werden. Jede dieser Functionen wird an der Stelle $P$ von einer bestimmten Ordnung

\origpage{409}
unendlich. Nach einem Satze des Herrn \emph{Weierstrass} kommen nun alle möglichen Ordnungen
\[
1,\, 2,\, 3,\, 4,\, \dots
\]
wirklich vor mit Ausnahme von $p$ Ordnungen. Diese $p$ fehlenden Ordnungen sind nun gerade die Zahlen $\varrho_1, \varrho_2, \dots \varrho_p${}^{*)}.

Für jede Stelle $P$, welche nicht mit einer der Stellen $P_1, P_2, \dots P_r$ zusammenfällt, ist $m = 0$ und folglich nach (4) und (5)
\[
\varrho_1 = 1,\, \varrho_2 = 2,\, \dots \varrho_p = p.
\]
Die niedrigste Ordnung, welche dann auftritt, ist also $p+1$. Umgekehrt, wenn $p+1$ die niedrigste auftretende Ordnung ist, so wird nothwendig $m = 0$.

Die Stellen $P_1, P_2, \dots P_\varrho$\ednote{Printed $P_\varrho$ instead of $P_r$; an apparent compositor misprint confusing italic r with Greek rho.} (für welche $m$ bez. die Werthe $m_1, m_2, \dots m_r$ besitzt) sind also vollständig dadurch charakterisirt, dass für jede derselben eine Function existirt, welche nur in ihr und zwar von geringerer als der $(p+1)^{\text{ten}}$ Ordnung unendlich wird. Hieraus erhellt, dass diese Stellen sich bei jeder eindeutigen Transformation der Fläche in sich nur unter einander vertauschen können${}^{**)}$.

Wir müssen nun die Unregelmässigkeiten, welche die Zahlenreihe $\varrho_1, \varrho_2, \dots \varrho_p$ darbieten kann, wenn $P$ mit einer der Stellen $P_1, P_2, \dots P_r$ zusammenfällt, näher untersuchen. Zunächst bemerken wir, dass $\varrho_p$ stets kleiner als $2p$ ist, weil $\frac{du_p}{dt}$ höchstens von der Ordnung $2p-2$ an der Stelle $P$ verschwinden kann. Es ist also
\[
\varrho_1 < \varrho_2 < \dots < \varrho_p \leqq 2p - 1.
\tag{6}
\]
Bezeichnen $z_\alpha$ und $z_\beta$ zwei Functionen, welche nur bei $P$ von den Ordnungen $\alpha$, bez. $\beta$ unendlich werden, so ist $z_\alpha z_\beta$ eine ähnliche Function mit der Ordnungszahl $\alpha + \beta$. Wenn daher $\alpha$ eine auftretende Ordnung, $\varrho$ eine fehlende bezeichnet, so muss auch $\varrho - \alpha$ eine fehlende Ordnung sein. Denn würde $\varrho - \alpha$ auftreten, so würde auch $(\varrho - \alpha) + \alpha = \varrho$ auftreten, gegen die Annahme.

Dies vorausgeschickt, sei $\alpha$ die \emph{kleinste} auftretende Ordnung. Es treten dann auch alle durch $\alpha$ theilbaren Ordnungen auf, woraus folgt, dass $\alpha \geqq 2$ ist.

Unter den Zahlen, welche $\equiv i \pmod \alpha$ sind, sei nun $r_i$ die grösste, welche unter den fehlenden Ordnungen vorkommt. Dann sind

\textbf{*)} Ich übergehe hier den Beweis, welcher übrigens leicht geführt wird durch Aufstellung aller jener Functionen vermöge der Integrale 2.~Gattung, welche bei $P$ unstetig werden. Man vgl. wegen des Satzes auch die Abhandlung \emph{Nöther}'s: „Beweis und Erweiterung eines algebraisch-functionentheoretischen Satzes des Herrn Weierstrass“. Crelle's Journal, Bd.~97.

\textbf{**)} Das Gleiche folgt auch schon aus der zweiten oben angegebenen Eigenschaft der Determinante $\Delta_u$.

\origpage{410}
alle fehlenden Ordnungen $\varrho_1, \varrho_2, \dots \varrho_p$ in folgender Tabelle enthalten:
\[
\left\{
\begin{aligned}
&1, & &1 + \alpha, & &1 + 2\alpha, \dots r_1, \\
&2, & &2 + \alpha, & &2 + 2\alpha, \dots r_2, \\
&\cdot & &\cdot & &\cdot \\
&\cdot & &\cdot & &\cdot \\
&\alpha - 1, & &(\alpha - 1) + \alpha, & &(\alpha - 1) + 2\alpha, \dots r_{\alpha-1},
\end{aligned}
\right.
\tag{7}
\]
wo in der ersten Horizontalreihe alle fehlenden Ordnungen stehen, welche $\equiv 1 \pmod \alpha$ sind u.~s.~w.${}^{*)}$ Da die Anzahl der Zahlen (7) gleich $p$ ist, so hat man
\[
p = \left(1 + \frac{r_1 - 1}{\alpha}\right) + \left(1 + \frac{r_2 - 2}{\alpha}\right) + \dots + \left(1 + \frac{r_{\alpha-1} - (\alpha - 1)}{\alpha}\right)
\]
oder
\[
r_1 + r_2 + \dots + r_{\alpha-1} = \alpha p - \frac{\alpha(\alpha - 1)}{2}.
\tag{8}
\]
Ferner ist die Summe der Zahlen (7) gleich
\[
\frac{1}{2}(r_1 + 1)\frac{r_1 + \alpha - 1}{\alpha} + \frac{1}{2}(r_2 + 2)\frac{r_2 + \alpha - 2}{\alpha} + \dots,
\]
und daher
\[
m = \frac{1}{2\alpha} \sum_1^{\alpha-1} (r_k + k)(r_k + \alpha - k) - \frac{1}{2} p(p+1).
\tag{9}
\]
Dieser Ausdruck lässt nun mit Hülfe von (8) die merkwürdige Umformung zu:
\[
m = \frac{p(p-1)}{2} - \frac{1}{2\alpha} \sum_1^{\alpha-1} r_k(2p - 1 - r_k) - \frac{(\alpha-1)(\alpha-2)}{6}.
\tag{10}
\]
Ist jetzt $\alpha = 2$, so folgt aus (8) und (10)
\[
m = \frac{p(p-1)}{2}.
\]
D.~h. „Ist die Fläche hyperelliptisch, so hat jeder der Punkte $P_1, \dots P_r$ die Multiplicität $\frac{p(p-1)}{2}$ und es ist also $r = 2p+2$.“

Diese Punkte sind keine anderen, wie die Nullstellen zweiter Ordnung von $dz$, wo $z$ die zweiwerthige Function der Fläche bedeutet.

\textbf{*)} Da $(r_i + \alpha) + (r_k + \alpha) = (r_i + r_k + \alpha) + \alpha$ eine auftretende Ordnung ist, so ist nothwendig $r_{i+k} \leqq r_i + r_k + \alpha$, wo die Indices (mod.~$\alpha$) zu reduciren sind. Diese Ungleichungen kommen jedoch nicht weiter in Betracht. Wegen der Aufstellung der fehlenden Ordnungen vgl. man übrigens „Schottky: „Ueber die conforme Abbildung mehrfach zusammenhängender ebener Flächen“ Crelle's Journal Bd.~83, pag.~308, 318, 319.

\origpage{411}
Wenn aber die Fläche nicht hyperelliptisch ist, so ist $\alpha$ nothwendig $> 2$ und die Gleichung (10) zeigt, dass
\[
m < \frac{p(p-1)}{2};
\]
denn $r_1, r_2, \dots r_{\alpha-1}$ sind positiv und
\[
2p - 1 - r_1,\, 2p - 1 - r_2,\, \dots 2p - 1 - r_{\alpha-1},
\]
nach (6) nicht negativ. Es sind also die Multiplicitäten $m_1, m_2, \dots m_r$ sämmtlich kleiner als $\frac{p(p-1)}{2}$, daher ihre Summe
\[
m_1 + m_2 + \dots + m_r = (p-1)p(p+1) < r \cdot \frac{p(p-1)}{2}
\]
und also
\[
r > 2p + 2.
\]

\emph{Die Anzahl der Stellen $P_1, P_2, \dots P_r$ übersteigt also auf jeder nicht hyperelliptischen Fläche stets die Zahl $2p+2$, womit auch die in Nr.~1 aufgestellte Behauptung gerechtfertigt ist.}

Die untere Grenze für die Zahl $r$ lässt sich durch eine eingehendere Discussion der Formel (10) noch erheblich vergrössern. Doch will ich hierauf nicht näher eingehen. Eine genaue Untersuchung würde übrigens nicht nur festzustellen haben, welche Werthe die Zahl $r$ annehmen kann, sondern auch welche Zahlsysteme $\varrho_1, \varrho_2, \dots \varrho_p$ für die einzelnen Stellen möglich sind. Bemerkt sei nur noch, dass im Falle $p = 3$ die untere Grenze der Zahl $r$ gleich $12$ ist, eine Thatsache, welche sich in folgender Weise aussprechen lässt:

„\emph{Eine ebene Curve 4.~Ordnung ohne vielfachen Punkt besitzt mindestens $12$ getrennt liegende Wendepunkte.}“

Das Beispiel der Curve
\[
f \equiv x_1^4 + x_2^4 + x_3^4 = 0
\]
zeigt, dass die untere Grenze $12$ wirklich erreicht wird. Die Hesse'sche Form von $f$ ist nämlich, abgesehen von einem Zahlenfactor, gleich $x_1^2 x_2^2 x_3^2$ und die Curve hat also nur die $12$ Schnittpunkte, welche auf den Linien $x_1 = 0$, $x_2 = 0$, $x_3 = 0$ liegen, zu Wendepunkten.

\subsection*{3.}

Die in der vorigen Nummer untersuchten Punkte $P_1, \dots P_r$ bilden das einfachste Beispiel einer „invarianten“ Punktgruppe einer Riemann'schen Fläche, d.~h. einer solchen Punktgruppe, deren Definition von einer besondern Darstellungsform der Fläche unabhängig ist. Für den in Nr.~1 gegebenen Beweis ist jede solche Punktgruppe brauchbar, wenn sie nur mehr als $2p+2$ Punkte enthält.

\origpage{412}
Es liegt nun die Vermuthung nahe, dass es auf jeder Fläche, deren Geschlecht $p > 1$ ist, stets unendlich viele invariante Punktgruppen giebt, und dass man mit solchen Punktgruppen die ganze Fläche überdecken kann. Ist die Fläche hyperelliptisch, so ist es leicht die Vermuthung zu bestätigen. In der That sei $y^2 = f(z)$ die gewöhnlich benutzte Darstellungsform der Fläche, welcher die Einführung der zweiwerthigen Function $z$ zu Grunde liegt. Dann bilden die Nullstellen jeder Covariante von $f(z)$ eine invariante Punktgruppe, und man erkennt ohne Schwierigkeit, dass mit derartigen Gruppen die ganze Fläche überzogen werden kann.

Ist die Fläche nicht hyperelliptisch, so stellen wir sie durch eine Curve $2p-2^{\text{ter}}$ Ordnung im Raume von $p-1$ Dimensionen dar, indem wir die homogenen Coordinaten $\varphi_1 : \varphi_2 : \dots : \varphi_p$ proportional setzen den $p$ Differentialen eines vollständigen Systems von Integralen 1.~Gattung${}^{*)}$. Jede covariante Fläche dieser Curve schneidet die letztere in einer invarianten Punktgruppe. Aber es tritt hier die Schwierigkeit auf, dass gewisse covariante Flächen nur existiren, wenn die $3p-3$ Moduln nicht besonderen Gleichungen genügen. Um so bemerkenswerther ist es, dass man trotzdem unendlich viele invariante Punktgruppen angeben kann, welche stets, also unabhängig von den besonderen Werthen der $3p-3$ Moduln, existiren. Ich will hierauf um so lieber eingehen, als die Bedeutung der oben betrachteten Punkte $P_1, P_2, \dots P_r$ dabei klarer hervortritt.

Man suche auf der Curve $2p-2^{\text{ter}}$ Ordnung diejenigen Punkte, in welchen eine Fläche $k^{\text{ter}}$ Ordnung hyperosculirt. Ist $k = 1$, handelt es sich also um die hyperosculirenden Ebenen, so sind die gesuchten Punkte keine anderen, wie die Punkte $P_1, P_2, \dots P_r$ und die Zahl dieser Punkte beträgt, wenn jeder mit der zugehörigen Multiplicität gezählt wird, $(p-1)p(p+1)$. Ich behaupte nun, dass auch für jeden Werth von $k > 1$, die gesuchten Punkte stets in endlicher Anzahl vorhanden sind, wie es folgender Satz näher ausspricht:

„\emph{Die Anzahl der Punkte auf der Curve $2p-2^{\text{ter}}$ Ordnung, in welchen dieselbe von einer Fläche $k^{\text{ter}}$ Ordnugg\ednote{Printed Ordnugg instead of Ordnung; an apparent compositor misprint.} hyperosculirt wird, beträgt stets}
\[
(2k-1)^2 \cdot p(p-1)^2.\text{“}
\]
Es wird vorausgesetzt, dass $k > 1$ sei und dass jeder Punkt mit der ihm zugehörigen, weiterhin näher zu bezeichnenden Multiplicität gezählt wird.

Der Beweis dieses Satzes beruht auf einem wichtigen Satze aus

\textbf{*)} Vgl. \emph{Weber}: „Ueber gewisse in der Theorie der Abel'schen Functionen auftretende Ausnahmefälle“. Mathem. Ann. Bd.~13, \emph{Nöther}: „Ueber die invariante Darstellung algebraischer Functionen“, ebenda, Bd.~17.

\origpage{413}
der Theorie der algebraischen Functionen, welchen Herr \emph{Nöther} aufgestellt und bewiesen hat${}^{*)}$.

Nach diesem Satz giebt es nämlich unter den Formen $k^{\text{ter}}$ Ordnung der Grössen $\varphi_1, \varphi_2, \dots \varphi_p$ genau $\varkappa = (2k-1)(p-1)$ linear unabhängige. Indem wir mit $u$ ein beliebig gewähltes Integral 1.~Gattung bezeichnen und $\varphi_1 = \frac{du_1}{du}$, $\varphi_2 = \frac{du_2}{du}$, $\dots \varphi_p = \frac{du_p}{du}$ annehmen, gehen die $\varkappa$ linear unabhängigen Formen --- sie seien
\[
\psi_1,\, \psi_2,\, \dots \psi_\varkappa \text{ ---}
\]
in Functionen des Ortes auf der Riemann'schen Fläche über. Die Punkte, in welchen eine Fläche $k^{\text{ter}}$ Ordnung hyperosculirt, sind nun die Nullstellen der Determinante:
\[
\Delta_\varkappa =
\begin{vmatrix}
\psi_1, & \psi_2, & \dots & \psi_\varkappa \\
\frac{d\psi_1}{du}, & \frac{d\psi_2}{du}, & \dots & \frac{d\psi_\varkappa}{du} \\
\cdot & \cdot & \dots & \cdot \\
\cdot & \cdot & \dots & \cdot \\
\frac{d^{\varkappa-1}\psi_1}{du^{\varkappa-1}}, & \frac{d^{\varkappa-1}\psi_2}{du^{\varkappa-1}}, & \dots & \frac{d^{\varkappa-1}\psi_\varkappa}{du^{\varkappa-1}}
\end{vmatrix}
\]
Diese Determinante verschwindet nicht identisch, weil $\psi_1, \psi_2, \dots \psi_\varkappa$ linear unabhängig sind. Sie stellt eine eindeutige algebraische Function des Ortes vor, welche an den $2p - 2$ Nullstellen von $du$ und zwar an jeder von der Ordnung $\frac{(2k-1+\varkappa)\varkappa}{2}$ unendlich wird. Folglich wird $\Delta_\varkappa$ im Ganzen
\[
(2p-2)(2k-1+\varkappa)\frac{\varkappa}{2} = (2k-1)^2 p(p-1)^2\text{-Mal}
\]
Null. Hiermit ist unser Satz bewiesen und es ist ersichtlich, dass die Multiplicität jedes einzelnen Punktes durch die Ordnung des Verschwindens von $\Delta_\varkappa$ angezeigt wird.

In dem einfachsten Falle $p = 3$ besagt unser Satz:

„\emph{Auf einer ebenen Curve 4.~Ordnung giebt es $12 \cdot (2k-1)^2$ Punkte in welchen eine Curve $k^{\text{ter}}$ Ordnung hyperosculirt.}“

Für den niedrigsten Werth von $k$, nämlich für $k = 2$, erhalten wir $12 \cdot 9 = 108$ Punkte, in welchen eine Curve 2.~Ordnung hyperosculirt. Diese Punkte sind die $24$ Wendepunkte und die $84$ sextactischen Punkte.

Bezeichnen $k, l, m, n$ vier ganze Zahlen, $\Delta_\varkappa, \Delta_\lambda, \Delta_\mu, \Delta_\nu$ die zugehörigen Determinanten, so überzeugt man sich leicht, dass die Function
\[
\frac{\lambda \, d\lg\Delta_\varkappa - \varkappa \, d\lg\Delta_\lambda}{\mu \, d\lg\Delta_\nu - \nu \, d\lg\Delta_\mu}
\]

\textbf{*)} Siehe das vorige Citat.

\origpage{414}
absolut invariant ist, dass also auch jede Punktgruppe, in welcher diese Function einen gegebenen Werth annimmt, eine invariante Gruppe ist. Mit solchen Gruppen lässt sich daher die Riemann'sche Fläche vollständig überdecken. Es bleibt freilich zu zeigen, dass $k, l, m, n$ stets so wählbar sind, dass jener Quotient sich nicht auf eine Constante reducirt. Doch will ich diesen Nachweis hier nicht führen.

\section*{II. Abschnitt.}

\subsection*{4.}

Bei den weiteren Untersuchungen werde ich eine bestimmte canonische Zerschneidung der Riemann'schen Flächen zu Grunde legen, welche ich hier kurz besprechen will. Nach dem Vorgange \emph{Riemann}'s pflegt man die Zerschneidung durch $p$ Schnittpaare und $p-1$ diese Paare verbindenden Schnitte auszuführen. Indessen ist es zweckmässiger die letzteren $p-1$ Schnitte vollständig zu beseitigen, da sie bei den meisten Untersuchungen eine ganz unnöthige Complication herbeiführen. Die $p-1$ Schnitte kommen in Fortfall, wenn wir folgende Zerschneidung wählen (welche übrigens auch durch einen Grenzübergang aus der Riemann'schen Zerschneidung abgeleitet werden kann${}^{*)}$). Wir wählen auf der Fläche einen Punkt $O$ beliebig und legen $2p$ knotenlose in $O$ beginnende und endigende Schnitte $A_1, B_1, \dots A_p, B_p$, durch welche die Fläche $F$ in eine einfach zusammenhängende $F'$ übergeht. Die Begrenzung von $F'$ wird durch die $4p$ Ufer der Schnitte gebildet. Die Bezeichnung und Aufeinanderfolge der Schnitte sei so festgelegt, dass man ihnen bei einem positiven Umlaufe um $O$ in der Folge
\[
A_1^+ B_1^+ A_1^- B_1^- \dots A_p^+ B_p^+ A_p^- B_p^-
\]
begegnet. Dies soll bedeuten, dass man zuerst den Schnitt $A_1$ vom negativen zum positiven Ufer, sodann ebenso den Schnitt $B_1$, sodann

\rmfigure{figures/fig-1.png}{Fig.~1.}{Canonical dissection of a Riemann surface of genus $p = 2$ with cuts radiating from a single point $O$.}

\textbf{*)} Vgl. \emph{F.~Klein}: „Neue Beiträge zur Riemann'schen Functionentheorie“. Mathem. Ann. Bd.~21, pag.~183\,ff. Man vergl. auch meine Arbeit: „Ueber Riemann'sche Flächen mit gegebenen Verzweigungspunkten“. Mathem. Annalen Bd.~39, pag.~1\,ff.

\origpage{415}
wieder $A_1$ vom positiven zum negativen Ufer überschreitet u.~s.~w. Die Fläche bietet also nach der Zerschneidung beim Punkte $O$ das durch Figur 1 (pag.~12) veranschaulichte Aussehen. Die Figur bezieht sich auf den Fall $p = 2$.

Es seien nun weiter $a_1, a_2, \dots a_w$ irgend $w$ Punkte der Fläche $F'$. Wir legen dann von $O$ aus durch das Innere von $F'$ knotenlose einander nicht treffende Schnitte $l_1, l_2, \dots l_w$ nach den Punkten $a_1, a_2, \dots a_w$. Die hierdurch aus $F'$ entstehende Fläche nennen wir $F''$ und wir nehmen die Aufeinanderfolge der Schnitte $l$ und ihre

\rmfigure{figures/fig-2.png}{Fig.~2.}{Canonical dissection of a Riemann surface of genus $p = 2$ with branch cuts radiating to points $a_1, a_2$.}

Lage gegen die Schnitte $A, B$ so an, dass ein positiver Umlauf um $O$ die Reihenfolge
\[
l_1^+ \, l_2^+ \, \dots \, l_w^+ \, A_1^+ \, B_1^+ \, A_1^- \, B_1^- \dots A_p^+ \, B_p^+ \, A_p^- \, B_p^-
\]
ergiebt. Man vergleiche Figur 2, welche den Fall $w = 2$, $p = 2$ veranschaulicht${}^{*)}$).

\subsection*{5.}

Man denke sich jetzt auf der Fläche $F$ eine analytische Function $Z$ des Ortes, welche keine natürlichen Grenzen und nur die Punkte $a_1, a_2, \dots a_w$ als Verzweigungspunkte besitzt. Der einzelne Zweig der Function $Z$ ist dann auf der Fläche $F''$ eindeutig. Ist

\textbf{*)} In den Figuren habe ich, dem Vorgange C.~Neumann's (Theorie der Abel'schen Functionen, Leipzig 1884) folgend, die positiven Ufer durch stärker gezeichnete Linien kenntlich gemacht. Die Pfeile geben den positiven Durchlaufungssinn für die Begrenzung der zerschnittenen Fläche an.

\origpage{416}
insbesondere $Z$ eine algebraische Function des Ortes, welche an der einzelnen Stelle der Fläche $F$ $n$ Werthe besitzt, so erscheint der ganze Werthevorrath der Function $Z$ auf der Fläche $F''$ in $n$ eindeutige Zweige zerlegt.

Wir nehmen dementsprechend $n$ in einander liegende Exemplare der Fläche $F''$, auf welche wir die $n$ Zweige der Function $Z$ ausbreiten und verbinden nun die $n$ Exemplare längs der Schnitte $l, A, B$ zu einer einzigen geschlossenen Fläche $\mathrm{F}$, auf welcher die Function $Z$ eindeutig ausgebreitet ist. Welches ist das Geschlecht $\mathrm{P}$ der so entstehenden Fläche $\mathrm{F}$? Wir wollen die $n$ Exemplare $F''$ als die „Blätter“ von $\mathrm{F}$ bezeichnen. Hängen dann an der Stelle $a_i$ die Blätter in $c_i$ Cyklen von je $v_1, v_2, \dots$ Blättern zusammen, so heisse
\[
(v_1 - 1) + (v_2 - 1) + \dots = n - c_i,
\]
wie üblich, die Multiplicität von $a_i$ als Verzweigungspunkt, und
\[
W = (n - c_1) + (n - c_2) + \dots + (n - c_w)
\]
die Gesammtzahl der Verzweigungen. Die aufgeworfene Frage wird dann durch nachstehende Gleichung beantwortet:
\[
2P - 2 = W + n(2p - 2).
\tag{1}
\]

Ist $p = 0$, so geht diese Gleichung in die bekannte Formel über, welche das Geschlecht einer $n$-blättrig über der complexen Zahlenebene ausgebreiteten Fläche berechnen lehrt. Ist $p > 0$, so kann man die Gleichung folgendermassen beweisen. Es sei $u$ ein überallendliches Integral der Fläche $F$, so ist $u$ ein ebensolches Integral für die Fläche $\mathrm{F}$. Es muss daher $du$ auf der Fläche $\mathrm{F}$ genau $2P - 2$ Mal von der zweiten Ordnung verschwinden. Diese Nullstellen von $du$ lassen sich aber sofort angeben. Es sind erstens die Verzweigungsstellen und zwar zählt die Verzweigungsstelle $a_i$ für $n - c_i$ Nullstellen, so dass wir, diesen Stellen entsprechend, im Ganzen $W$ Nullstellen erhalten. Sodann verschwindet zweitens $du$ auf der Fläche $F$ an $2p - 2$ Stellen von der zweiten Ordnung und diese Stellen ergeben auf $\mathrm{F}$ im Ganzen $n(2p - 2)$ Nullstellen, da jede Stelle von $F$ sich in $n$ Exemplaren auf der Fläche $\mathrm{F}$ wiederfindet. Es ist also $2P - 2 = W + n(2p - 2)$, w.~z.~b.~w.${}^{*)}$

Ich will hier beiläufig aus der Gleichung (1) eine Folgerung ziehen. Es seien zwei Riemann'sche Flächen $F_1$ und $F_2$ vom Geschlecht $p_1$ bez. $p_2$ algebraisch so auf einander bezogen, dass jeder Stelle von $F_1$ $n_1$ Stellen von $F_2$ und umgekehrt jeder Stelle von $F_2$ $n_2$ Stellen von $F_1$ correspondiren. Diese Beziehung können wir dadurch zu einer

\textbf{*)} Ein anderer auf Betrachtungen der Analysis situs beruhender Beweis findet sich in meiner Arbeit „Ueber Riemann'sche Flächen mit gegebenen Verzweigungspunkten.“

\origpage{417}
eindeutigen machen, dass wir die Fläche $F_1$ $n_1$-fach, die Fläche $F_2$ $n_2$-fach überdecken. Da die so entstehenden über $F_1$ und $F_2$ ausgebreiteten Flächen gleiches Geschlecht $P$ haben müssen, so ist nach (1)
\[
W_1 + n_1(2p_1 - 2) = W_2 + n_2(2p_2 - 2).
\tag{2}
\]
Hier bedeutet $W_1$ die Anzahl der Stellen $P$ auf $F_1$, welchen zusammenfallende Stellen auf $F_2$ entsprechen, und zwar so, dass einem Umlauf um $P$ eine Vertauschung der entsprechenden Stellen auf $P_2$\ednote{Printed $P_2$ instead of $F_2$; an apparent compositor misprint.} entspricht. Eine analoge Bedeutung besitzt $W_2$. Man erkennt in (2) die von Herrn \emph{Zeuthen} gegebene Relation zwischen den Geschlechtszahlen zweier mehrdeutig auf einander bezogener algebraischen Gebilde.${}^{*)}$ Unsere Ableitung der Relation zeigt zugleich unzweideutig, mit welchen Multiplicitäten diejenigen Punkte des einzelnen Gebildes zu zählen sind, denen auf dem anderen Gebilde zusammenfallende Punkte entsprechen.

\subsection*{6.}

Ich wende mich jetzt zu unserem eigentlichen Gegenstande zurück. Es sei $F$ eine Riemann'sche Fläche von beliebigem Geschlecht. (Die Fälle $p = 0$ und $p = 1$ werden also nicht ausgeschlossen.) Wir nehmen an, dass diese Fläche eine Gruppe von $r$ eindeutigen Transformationen in sich besitze, die wir mit
\[
S_1, S_2, \dots S_r
\tag{1}
\]
bezeichnen. Die aus $S_i$ und $S_k$ zusammengesetzte Transformation $S_i S_k$ wollen wir auch mit $S_{ik}$ bezeichnen, so dass $ik$ wieder eine der Zahlen $1, 2, \dots r$ bedeutet. $S_1$ sei die identische Transformation, also $S_{1i} = S_{i1} = S_i$. Die Stellen der Fläche ordnen sich nun in Gruppen zu je $r$ an, wie
\[
P_1, P_2, \dots P_r,
\tag{2}
\]
welche dadurch erhalten werden, dass man eine willkürlich gewählte Stelle $P_1$ mit denjenigen Stellen zusammenfasst, welche aus ihr durch die Transformationen (1) hervorgehen. Die Gruppe (2) wird aus $r$ getrennt liegenden Punkten bestehen, wenn $P_1$ bei keiner der Transformationen (1), ausser der Identität, festbleibt. Wenn jedoch $P_1$ bei $k$ Transformationen (1) fest bleibt, so bilden diese für sich eine Gruppe und man erkennt hieraus leicht, dass in diesem Falle aus $P_1$ nur $\frac{r}{k}$ verschiedene Punkte hervorgehen, welche je $k$ Punkte in sich vereinigen.

Wir betrachten nun die Classe der eindeutiger\ednote{Printed eindeutiger instead of eindeutigen; an apparent compositor error.} algebraischen Functionen der Fläche $F$. Ist $z$ eine dieser Functionen und bedeuten

\textbf{*)} Mathem. Ann. Bd.~3, p.~150.

\origpage{418}
$z_1, z_2, \dots z_r$ die Werthe von $z$ an den Stellen $P_1, P_2, \dots P_r$, so können wir die Summe
\[
Z = z_1 + z_2 + \dots + z_r
\]
als eine Function der Stelle $P_1$ ansehen. Als solche hat sie die Eigenschaften, erstens eine eindeutige algebraische Function auf $F$ zu sein (also der betrachteten Classe anzugehören) und zweitens der Bedingung
\[
Z_1 = Z_2 = \dots = Z_r
\]
zu genügen, wo $Z_1, Z_2, \dots Z_r$ die Werthe von $Z$ in den Punkten irgend einer Gruppe bezeichnen. Da man noch über die Unendlichkeitspunkte von $z$ verfügen kann, so leuchtet ein, dass es unendlich viele Functionen $Z$ von diesen beiden Eigenschaften giebt. Die Gesammtheit der Functionen $Z$ bildet wieder eine Classe von algebraischen Functionen, da Summe, Differenz, Product und Quotient zweier Functionen $Z$ wieder eine ebensolche Function darstellt. Die Riemann'sche Fläche, auf welcher die Functionen $Z$ die Gesammtheit der eindeutigen algebraischen Functionen des Ortes bilden, heisse $\Phi$.${}^{*)}$

Die Fläche $\Phi$ ist nun derartig auf die Fläche $F$ bezogen, dass jeder Stelle von $F$ eine einzige Stelle von $\Phi$, dagegen umgekehrt jeder Stelle von $\Phi$ $r$ Stellen von $F$, nämlich die $r$ Stellen einer Gruppe (2) entsprechen. Diejenigen Stellen der Fläche $\Phi$, welchen Gruppen von weniger als $r$ (also Gruppen von $\frac{r}{k}$) Stellen auf der Fläche $F$ entsprechen, wollen wir mit
\[
a_1, a_2, \dots a_w
\tag{3}
\]
bezeichnen. Wenn wir nun eine Stelle $P$, welche nicht zu den Stellen (3) gehört, sich auf der Fläche $\Phi$ in Bewegung setzen und einen die Stellen (3) vermeidenden geschlossenen Weg durchlaufen lassen, so werden die entsprechenden Stellen $P_1, P_2, \dots P_r$ auf der Fläche $F$ sich zugleich bewegen und es können sich schliesslich, wenn $P$ die alte Lage wieder erreicht hat, $P_1, P_2, \dots P_r$ nur untereinander vertauscht haben. Wenn hierbei $P_1$ in $P_i$ übergeht, so geht $P_k$ nothwendig in $P_{ik}$ über. Denn um den Weg von $P_k$ zu verfolgen, brauchen wir nur auf jeden Punkt des Weges von $P_1$ die Transformation $S_k$ anzuwenden; durch diese Transformation geht aber $P_i$ in $P_{ik}$ über. Es erfahren also die Stellen $P_1, P_2, \dots P_r$, wenn $P$ einen geschlossenen Weg durchläuft, eine Vertauschung der Gestalt
\[
\begin{pmatrix}
P_1, & P_2, & \dots & P_r \\
P_{i1}, & P_{i2}, & \dots & P_{ir}
\end{pmatrix}.
\tag{4}
\]

\textbf{*)} Ist $Z_0$ eine bestimmte der Functionen $Z$, so kann man für $\Phi$ z.~B. diejenige über der $Z_0$-Ebene ausgebreitete Fläche nehmen, welche die Verzweigung der Functionen $Z$ in Bezug auf $Z_0$ darstellt.

\origpage{419}
Dies vorausgeschickt, möge jetzt $\Phi$ in der oben geschilderten Weise in die Fläche $\Phi''$ zerschnitten werden. Ist $\pi$ das Geschlecht von $\Phi$, so geschieht die Zerschneidung durch $w + 2\pi$ Schnitte, welche wir (wie oben) mit
\[
l_1, \ l_2, \ \dots \ l_w, \ A_1, \ B_1, \ \dots \ A_\pi, \ B_\pi
\tag{5}
\]
bezeichnen. Wir fixiren im Innern von $\Phi''$ einen Punkt $P'$, dem auf der Fläche $F$ die Punkte $P_1', P_2', \dots P_r'$ entsprechen mögen; ferner nehmen wir $r$ übereinander liegende Exemplare (Blätter) $\Phi''$ an, bezeichnen dieselben in irgend einer Reihenfolge mit denselben Buchstaben, wie die Transformationen (1), also mit
\[
S_1, \ S_2, \ \dots \ S_r
\]
und ordnen dem Punkte $P'$ des Blattes $S_k$ den Punkt $P_k'$ auf der Fläche $F$ zu. Bewegt sich $P'$ auf dem Blatte $S_k$ nach $P$, so wird sich der Punkt $P_k'$ auf $F$ nach $P_k$ bewegen und zwar wird die Endlage $P_k$ dieselbe sein, auf welchem Wege auch $P'$ nach $P$ geführt wird. Für jede Stelle $P$ erscheinen also die entsprechenden Punkte $P_1, P_2, \dots P_r$ in bestimmter Weise den $r$ Blättern $S_1, S_2, \dots S_r$ zugeordnet. Wenn wir jetzt die Blätter $S_1, S_2, \dots S_r$ längs der Schnitte (5) in geeigneter Weise verbinden, so erhalten wir eine über $\Phi$ $r$-blättrig ausgebreitete Fläche, welche eindeutig umkehrbar auf $F$ bezogen ist und also mit $F$ in dem Sinne identisch ist, als sie gerade so gut wie $F$ als Träger der durch $F$ dargestellten Classe von algebraischen Functionen angesehen werden kann. Was nun die Verbindung der Blätter angeht, so beachte man, dass bei einem von $P$ beschriebenen geschlossenen Wege $P_1, P_2, \dots P_r$ eine Vertauschung der Gestalt (4) erfahren. Es müssen also z.~B. längs $l_1$ die Blätter so verbunden werden, dass man beim Uebertritt von dem negativen zum positiven Ufer von $l_1$ aus den Blättern
\[
S_1, \quad S_2, \quad \dots \quad S_r
\]
bez. in die Blätter
\[
S_i S_1, \quad S_i S_2, \quad \dots \quad S_i S_r
\]
gelangt.

Jedem der Schnitte (5) ist also eine bestimmte Transformation $S_i$ zugeordnet. Ist der Reihe nach für den Schnitt
\[
\begin{gathered}
l_1, \ l_2, \ \dots \ l_w, \ A_1, \ B_1, \ \dots \ A_\pi, \ B_\pi, \\
S_i = T_1, \ T_2, \ \dots \ T_w, \ U_1, \ V_1, \ \dots \ U_\pi, \ V_\pi,
\end{gathered}
\]
wo also die $T, U, V$ bestimmte der Transformationen (1) bezeichnen, so können wir die $r$-blättrig über $\Phi$ ausgebreitete Fläche $F$ schematisch durch
\[
F = \begin{pmatrix}
l_1, & l_2, & \dots & l_w, & A_1, & B_1, & \dots & A_\pi, & B_\pi \\
T_1, & T_2, & \dots & T_w, & U_1, & V_1, & \dots & U_\pi, & V_\pi
\end{pmatrix}
\tag{6}
\]

\origpage{420}
charakterisiren. Die Transformationen $T, U, V$ genügen den folgenden Bedingungen:

Erstens ist
\[
T_1 T_2 \dots T_w \, U_1 \, V_1 \, U_1^{-1} \, V_1^{-1} \dots U_\pi \, V_\pi \, U_\pi^{-1} \, V_\pi^{-1} = S_1 = 1.
\tag{7}
\]

Zweitens lässt sich die ganze Gruppe $S_1, S_2, \dots S_r$ durch Combination der $T, U, V$ erzeugen. Die erste Bedingung (7) sagt nichts Anderes aus, als dass der Punkt $O$, von welchem die Schnitte $l, A, B$ auslaufen, kein Verzweigungspunkt ist. Die zweite Bedingung besagt, dass die $r$-blättrige Fläche $F$ in sich zusammenhängt. Diese Betrachtung lehrt nun, dass man die allgemeinste Riemann'sche Fläche, welche eine Gruppe eindeutiger Transformationen in sich besitzt, auf folgendem Wege erhält:

\emph{Man bezeichne mit}
\[
S_1, \ S_2, \ \dots \ S_r
\]
\emph{irgend $r$ Operationen, welche eine Gruppe bilden. Ferner seien}
\[
T_1, \ T_2, \ \dots \ T_w, \quad U_1, \ V_1, \ \dots \ U_\pi, \ V_\pi
\]
\emph{irgend $w + 2\pi$ dieser Operationen, durch welche sich die ganze Gruppe erzeugen lässt und welche der Bedingung}
\[
T_1 T_2 \dots T_w \, U_1 \, V_1 \, U_1^{-1} \, V_1^{-1} \dots U_\pi \, V_\pi \, U_\pi^{-1} \, V_\pi^{-1} = 1
\]
\emph{genügen. Man nehme eine beliebige Riemann'sche Fläche $\Phi$ vom Geschlecht $\pi$, markire auf dieser irgend $w$ Punkte $a_1, a_2, \dots a_w$ und zerschneide dieselbe in der oben (Nr.~4) geschilderten Weise durch Schnitte $l_1, l_2, \dots l_w, A_1, B_1, \dots A_\pi, B_\pi$ in eine einfach zusammenhängende Fläche $\Phi''$. Endlich decke man $r$ Exemplare der Fläche $\Phi''$, welche man in irgend einer Reihenfolge mit}
\[
S_1, \ S_2, \ \dots \ S_r
\]
\emph{(also genau wie die gegebenen Operationen) bezeichnet, übereinander und verbinde dieselben längs der Schnitte $l, A, B$ zu einer einzigen geschlossenen Fläche, wie es das Schema}
\[
F = \begin{pmatrix}
l_1, & l_2, & \dots & l_w, & A_1, & B_1, & \dots & A_\pi, & B_\pi \\
T_1, & T_2, & \dots & T_w, & U_1, & V_1, & \dots & U_\pi, & V_\pi
\end{pmatrix}
\]
\emph{andeutet. Die auf diese Weise entstehenden Flächen umfassen alle Riemann'sche Flächen, welche eine Gruppe von eindeutigen Transformationen in sich besitzen.}

Hierbei ist zu bemerken, dass das für $F$ gegebene Schema folgende Bedeutung hat: An dem Schnitte $l_1$ sollen die $r$ Blätter $\Phi''$ so verbunden werden, dass man beim Uebertritt von der negativen auf die positive Seite des Schnittes aus den Blättern
\[
S_1, \quad S_2, \quad \dots \quad S_r
\]
bez. in die Blätter
\[
T_1 S_1, \quad T_1 S_2, \quad \dots \quad T_1 S_r
\]

\origpage{421}
gelangt. Entsprechendes ist für die Schnitte $l_2, \dots l_w, A_1, B_1, \dots A_\pi, B_\pi$ zu bemerken. Wenn nun feststeht, dass durch die obige Construction alle Flächen mit einer Gruppe von $r$ eindeutigen Transformationen in sich erhalten werden, so erhebt sich jetzt umgekehrt die Frage, ob auch \emph{jede} durch die Construction hergestellte Fläche eine Gruppe von $r$ Transformationen in sich besitzt. Dass diese Frage zu bejahen ist, ergiebt sich aus folgender Betrachtung. Man fixire auf der $r$-blättrigen nach obiger Vorschrift hergestellten Fläche $F$ irgend $r$ übereinanderliegende Punkte und bezeichne sie wie die Blätter, in welchen sie liegen. Ordnet man nun den Punkten
\[
S_1, \quad S_2, \quad \dots \quad S_r
\]
bez. die Punkte
\[
S_1 S_i, \quad S_2 S_i, \quad \dots \quad S_r S_i
\]
zu, wo $S_i$ irgend eine der $r$ Operationen bedeutet, so ist hierdurch eine eindeutige Transformation von $F$ in sich bestimmt. Denn bei Fortsetzung dieser Zuordnung könnte eine Mehrdeutigkeit nur durch Ueberschreitung eines Schnittes entstehen. Eine solche Mehrdeutigkeit entsteht aber thatsächlich nicht, weil bei der Ueberschreitung von $l_1$ z.~B. die ursprüngliche Zuordnung in
\[
\begin{gathered}
T_1 S_1, \quad T_1 S_2, \quad \dots \quad T_1 S_r, \\
T_1 S_1 S_i, \quad T_1 S_2 S_i, \quad \dots \quad T_1 S_r S_i
\end{gathered}
\]
übergeht, welche ersichtlich mit der ursprünglichen identisch ist. Den $r$ Operationen $S_i$ entsprechend, besitzt also die Fläche $r$ eindeutige Transformationen in sich. Da die Gruppe der $r$ Operationen $S_1, S_2, \dots S_r$ ebenso wie die Fläche $\Phi$ ganz willkürlich angenommen werden kann und die Operationen $T, U, V$ dann noch in mannigfaltiger Weise gewählt werden können, so ergiebt sich beiläufig der Satz:

\emph{Ist eine endliche Gruppe gegeben, so kann man auf mannigfaltige Weise eine Riemann'sche Fläche herstellen, welche eine zu der gegebenen Gruppe holoedrisch isomorphe Gruppe von eindeutigen Transformationen in sich besitzt.}${}^{*)}$

Unsere Betrachtungen zeigen zugleich, wie man die allgemeinste derartige Fläche findet.

\subsection*{7.}

Wir betrachten jetzt einc\ednote{Printed einc instead of eine; an apparent compositor misprint.} Riemann'sche Fläche $F$, welche nach den Angaben der vorigen Nummer $r$-blättrig über einer Fläche $\Phi$

\textbf{*)} Die an die Untersuchungen von F.~Klein (Math. Ann. Bde.~9--17) anknüpfenden Arbeiten von W.~Dyck (Math. Ann, Bde.~17 und 20) über reguläre Riemann'sche Flächen bieten mannigfaltige Berührungspunkte mit den obigen Betrachtungen. Was den zuletzt angeführten Satz angeht, so vergleiche man den Satz von Dyck, Math. Ann. Bd.~20, p.~30.

\origpage{422}
vom Geschlecht $\pi$ ausgebreitet ist. Wir wissen, dass $F$ die allgemeinste Fläche mit einer Gruppe von $r$ eindeutigen Transformationen in sich darstellt. Es sei unsere nächste Aufgabe, das Geschlecht von $F$ zu berechnen. Bei einer Umkreisung des Punktes $a_1$ geht nun allgemein das Blatt $S_i$ in das Blatt $T_1 S_i$ über. Ist daher $k_1$ die Ordnung der Transformation $T_1$, so hängen die Blätter bei $a_1$ in Cyklen von je $k_1$ Blättern zusammen, wie
\[
(S_i, \ T_1 S_i, \ T_1^2 S_i, \ \dots \ T_1^{k_1-1} S_i).
\]
Der Punkt $a_1$ hat also als Verzweigungspunkt die Multiplicität
\[
r - \frac{r}{k_1} = r \Bigl(1 - \frac{1}{k_1}\Bigr)
\]
und ebenso haben $a_2, a_3, \dots a_w$ bez. die Multiplicitäten
\[
r \Bigl(1 - \frac{1}{k_2}\Bigr), \quad r \Bigl(1 - \frac{1}{k_3}\Bigr), \quad \dots \quad r \Bigl(1 - \frac{1}{k_w}\Bigr),
\]
wenn $k_2, k_3, \dots k_w$ die Ordnungen von $T_2, T_3, \dots T_w$ bedeuten.

Nach Nr.~5 ist daher das Geschlecht $p$ der Fläche $F$ aus der Gleichung
\[
\frac{2p-2}{r} = 2\pi - 2 + \Bigl(1 - \frac{1}{k_1}\Bigr) + \Bigl(1 - \frac{1}{k_2}\Bigr) + \dots + \Bigl(1 - \frac{1}{k_w}\Bigr)
\tag{1}
\]
zu berechnen. Wir wollen jetzt annehmen, dass die Fläche $F$ von einem Geschlecht $p > 1$ sei, und unter dieser Annahme die Gleichung (1) näher betrachten. Wir unterscheiden dabei drei Fälle, je nachdem $\pi \geqq 2$, $\pi = 1$, $\pi = 0$ ist.

Im ersten Falle $\pi \geqq 2$ lehrt die Gleichung (1), dass $\frac{2p-2}{r} \geqq 2$, also
\[
r \leqq p - 1
\tag{2}
\]
ist. Im zweiten Falle $\pi = 1$ kann $w$ nicht $= 0$ sein, weil sonst nach Gleichung (1) auch $p = 1$ sein würde. Es ist daher
\[
\frac{2p-2}{r} \geqq \Bigl(1 - \frac{1}{k_1}\Bigr) \geqq \frac{1}{2},
\]
weil $k_1$ mindestens gleich $2$ ist. Es kommt also jetzt
\[
r \leqq 4(p - 1).
\tag{3}
\]

Im dritten Falle $\pi = 0$ haben wir nach (1)
\[
\begin{aligned}
\frac{2p-2}{r} &= \Bigl(1 - \frac{1}{k_1}\Bigr) + \Bigl(1 - \frac{1}{k_2}\Bigr) + \dots + \Bigl(1 - \frac{1}{k_w}\Bigr) - 2 \\
&= w - 2 - \frac{1}{k_1} - \frac{1}{k_2} - \dots - \frac{1}{k_w};
\end{aligned}
\tag{1'}
\]
es muss daher $w$ mindestens gleich $3$ sein. Wir unterscheiden nun drei Unterfälle. Ist erstens $w \geqq 5$, so erhalten wir aus ($1'$)

\origpage{423}
\[
\frac{2p - 2}{r} \geqq \frac{w}{2} - 2 \geqq \frac{5}{2} - 2 = \frac{1}{2},
\]
also
\[
r \leqq 4(p - 1).
\tag{4}
\]

Ist zweitens $w = 4$, also
\[
\frac{2p - 2}{r} = 2 - \frac{1}{k_1} - \frac{1}{k_2} - \frac{1}{k_3} - \frac{1}{k_4},
\]
so haben wir, indem wir $k_1 \leqq k_2 \leqq k_3 \leqq k_4$ annehmen, folgende tabellarisch angeordnete Möglichkeiten:
\[
\begin{array}{c|c|c|c|c|c}
\hline\hline
k_1 & k_2 & k_3 & k_4 & \frac{2p - 2}{r} & \\
\hline
> 2 & > 2 & > 2 & > 2 & \geqq \frac{2}{3} & r \leqq 3(p - 1) \\
\hline
= 2 & > 2 & > 2 & > 2 & \geqq \frac{1}{2} & r \leqq 4(p - 1) \\
\hline
= 2 & = 2 & > 2 & > 2 & \geqq \frac{1}{3} & r \leqq 6(p - 1) \\
\hline
= 2 & = 2 & = 2 & > 2 & \geqq \frac{1}{6} & r \leqq 12(p - 1)
\end{array}
\tag{5}
\]

Ist endlich drittens $w = 3$, also
\[
\frac{2p - 2}{r} = 1 - \frac{1}{k_1} - \frac{1}{k_2} - \frac{1}{k_3},
\]
so haben wir, indem wir $k_1 \leqq k_2 \leqq k_3$ annehmen, folgende wiederum tabellarisch aufgeführte Möglichkeiten:
\[
\begin{array}{c|c|c|c|c}
\hline\hline
k_1 & k_2 & k_3 & \frac{2p - 2}{r} & \\
\hline
> 3 & > 3 & > 3 & \geqq \frac{1}{4} & r \leqq 8(p - 1) \\
\hline
= 3 & > 3 & > 3 & \geqq \frac{1}{6} & r \leqq 12(p - 1) \\
\hline
= 3 & 3 & > 3 & \geqq \frac{1}{12} & r \leqq 24(p - 1) \\
\hline
= 2 & > 4 & > 4 & \geqq \frac{1}{10} & r \leqq 20(p - 1) \\
\hline
= 2 & = 4 & > 4 & \geqq \frac{1}{20} & r \leqq 40(p - 1) \\
\hline
= 2 & = 3 & > 6 & \geqq \frac{1}{42} & r \leqq 84(p - 1).
\end{array}
\tag{6}
\]

Ein Blick auf die Ungleichungen (2), (3), (4), (5), (6)\ednote{In Tabelle (6) druckte der Setzer in der dritten Zeile 3 statt = 3; in der sechsten Zeile ist die 6 mit Tinte über eine gedruckte 5 gebessert.} zeigt, dass $r$ die Zahl $84(p - 1)$ nicht übersteigen kann. Mit anderen Worten, es besteht der Satz:

\origpage{424}
„\emph{Die Anzahl der eindeutigen Transformationen in sich, welche eine Riemann'sche Fläche vom Geschlecht $p > 1$ besitzen kann, beträgt im Maximum $84(p - 1)$.}“

Zugleich lehrt unsere Discussion, dass das Maximum von $84(p - 1)$ Transformationen nur für solche Flächen eintreten kann, welche über einer Fläche vom Geschlecht Null oder, was auf dasselbe hinauskommt, über der complexen Zahlenebene so ausgebreitet sind, dass die Blätter an einer Stelle zu je 2, an einer anderen Stelle zu je 3 und an einer dritten Stelle zu je 7 im Cyklus zusammenhängen. (Der letzte Fall der Tabelle (6)). Diese Flächen hängen ersichtlich auf's Engste mit der Theorie derjenigen Schwarz'schen $s$-Function zusammen, welche die Halbebene auf ein Kreisbogendreieck mit den Winkeln $\frac{\pi}{2}$, $\frac{\pi}{3}$, $\frac{\pi}{7}$, abbildet.*)

Für welche Werthe von $p$ das Maximum $84(p - 1)$ wirklich erreicht wird, habe ich nicht untersucht. Bemerkt sei nur noch, dass für $p = 3$ eine Fläche mit $84(p - 1) = 168$ Transformationen in sich existirt; es ist die Fläche, welche zur Galois'schen Resolvente der Modulargleichung $7^{\text{ter}}$ Ordnung gehört, die zuerst von Herrn Klein betrachtet und ausführlich untersucht worden ist.**)

\subsection*{8.}

Zur Erläuterung der in Nr.~6 gegebenen allgemeinen Construction der Flächen mit eindeutigen Transformationen in sich mögen hier noch einige Beispiele folgen. Es sei zuerst
\[
S_1, \ S_2, \ \dots \ S_{60}
\tag{1}
\]
das System der geraden Vertauschungen von fünf Dingen $x_1, x_2, x_3, x_4, x_5$. Greifen wir aus dem System (1) die folgenden Vertauschungen heraus
\[
T_1 = (x_1 x_2)(x_3 x_4), \quad T_2 = (x_2 x_5 x_4), \quad T_3 = (x_1 x_2 x_3 x_4 x_5),
\]
so lässt sich aus diesen die ganze Gruppe (1) erzeugen, und es ist
\[
T_1 T_2 T_3 = 1.
\]
Nehmen wir nun eine Riemann'sche Fläche vom Geschlecht Null, also am einfachsten die complexe Zahlenebene, in dieser drei beliebige Punkte $a_1$, $a_2$, $a_3$, die wir mit einem vierten Punkt $O$ durch die Schnitte $l_1$, $l_2$, $l_3$ verbinden! Von der zerschnittenen Ebene decken wir 60 Exemplare, welche wir in irgend einer Reihenfolge mit

\textbf{*)} Vgl. Klein-Fricke: „Vorlesungen über die Theorie der elliptischen Modulfunctionen“ (Leipzig, 1890) die Figur pag.~109.

\textbf{**)} Mathematische Annalen Bd.~14, pag.~428. Siehe auch Capitel 6 des im vorigen Citat genannten Werkes.

\origpage{425}
$S_1, S_2, \dots S_{60}$ bezeichnen, übereinander und verbinden dieselben zu der Fläche
\[
F = \begin{pmatrix} l_1, & l_2, & l_3 \\ T_1, & T_2, & T_3 \end{pmatrix}.
\]
Die Verbindung geschieht nach der Massgabe, dass man beim Uebertritt vom negativen zum positiven Ufer der Linie $l_i$ aus den Blättern
\[
S_1, \quad S_2, \quad \dots \quad S_{60}
\]
bez. in die Blätter
\[
T_i S_1, \quad T_i S_2, \quad \dots \quad T_i S_{60}
\]
gelangen soll. Die so entstehende Fläche $F$ besitzt eine Gruppe von 60 eindeutigen Transformationen in sich, welche zu der Gruppe (1) der geraden Vertauschungen holoedrisch-isomorph ist. Das Geschlecht $p$ von $F$ ergiebt sich aus Gleichung (1) in Nr.~7, indem man
\[
r = 60, \quad \pi = 0, \quad k_1 = 2, \quad k_2 = 3, \quad k_3 = 5
\]
setzt. Man erhält
\[
\frac{2p - 2}{60} = - 2 + \frac{1}{2} + \frac{2}{3} + \frac{4}{5} = - \frac{1}{30},
\]
also $p = 0$, ein Resultat, an welches sich die von Herrn Klein gegebene Theorie der Gleichungen fünften Grades anschliessen lässt.*)

Als ein anderes Beispiel nehmen wir das System
\[
S_1, \ S_2, \ \dots \ S_{n!}
\tag{2}
\]
aller Vertauschungen von $n$ Dingen $x_1, x_2, \dots x_n$. Wir wählen aus (2)
\[
T_1 = (x_1 x_2), \quad T_2 = (x_n x_{n-1} \dots x_2), \quad T_3 = (x_1 x_2 \dots x_n),
\]
welche die Bedingung $T_1 T_2 T_3 = 1$ erfüllen und ein System erzeugender Substitutionen bilden. Indem wir in der complexen Zahlenebene von einem Punkte $O$ aus die Schnitte $l_1$, $l_2$, $l_3$ nach irgend drei Punkten $a_1$, $a_2$, $a_3$ legen, $n!$ Exemplare der zerschnittenen Ebene übereinanderdecken und endlich diese Exemplare zur Fläche
\[
F = \begin{pmatrix} l_1 & l_2 & l_3 \\ T_1 & T_2 & T_3 \end{pmatrix}
\]
vereinigen, erhalten wir eine Fläche mit $n!$ eindeutigen Transformationen in sich. Die letzteren bilden eine zur Gruppe der Vertauschungen von $n$ Dingen isomorphe Gruppe. Für das Geschlecht $p$ der Fläche $F$ finden wir nach Nr.~7 (da $\pi = 0$, $k_1 = 2$, $k_2 = n - 1$, $k_3 = n$ ist) den Werth

\textbf{*)} F.~Klein: „Vorlesungen über das Ikosaeder und die Auflösung der Gleichungen vom fünften Grade.“ Leipzig 1884.

\origpage{426}
\[
p = 1 + \frac{1}{4}\,(n - 2)!\,(n^2 - 5n + 2)
\tag{3}
\]
also für $n = 3, 4, 5, 6 \dots$ der Reihe nach $p = 0, 0, 4, 49, \dots$.

Die Fläche $F$ gestattet eine sehr einfache analytische Darstellung. Man betrachte die Gleichung
\[
s^n - \frac{n}{n - 1}\,s^{n-1} = z.
\tag{4}
\]
Die Verzweigung von $s$ in Bezug auf $z$ wird durch eine $n$-blättrig über der $z$-Ebene ausgebreitete Fläche dargestellt. Die Blätter der Fläche hängen bei $z = \infty$ im Cyklus zusammen, bei $z = 0$ hängen $n - 1$ Blätter im Cyklus zusammen, eines verläuft isolirt, endlich bei $z = \frac{-1}{n-1}$ findet eine einfache Verzweigung statt. Offenbar stellt daher (für $a_1 = -\frac{1}{n-1}$, $a_2 = 0$, $a_3 = \infty$) unsere Fläche $F$ die Verzweigung der Galois'schen Resolvente von (4) dar, d.~h. die Verzweigung der simultan betrachteten Wurzeln $s_1, s_2, \dots s_n$ von (4) in Bezug auf $z$. Nun ist aber die einzelne Stelle $P$ der Fläche $F$, wie man sich sofort überzeugt, schon durch die Verhältnisse der Grössen $s_1, s_2, \dots s_n$ festgelegt. Lassen wir daher $P$ die Fläche $F$ durchlaufen, so wird der Punkt
\[
y_1 : y_2 : \dots : y_n = \frac{1}{s_1} : \frac{1}{s_2} : \dots : \frac{1}{s_n}
\]
eine eindeutig umkehrbar auf die Fläche $F$ bezogene Curve beschreiben, wenn wir $y_1, y_2, \dots y_n$ als homogene Coordinaten in einem Raume von $n - 1$ Dimensionen deuten. Aus der Gestalt der Gleichung (4) folgt aber, dass diese Curve durch die folgenden $n - 2$ Gleichungen definirt werden kann:
\[
\left\{
\begin{aligned}
y_1^{\phantom{n-2}} + y_2^{\phantom{n-2}} + \dots + y_n^{\phantom{n-2}} &= 0, \\
y_1^2 + y_2^2 + \dots + y_n^2 &= 0, \\
\dots\dots\dots\dots\dots\dots & \\
y_1^{n-2} + y_2^{n-2} + \dots + y_n^{n-2} &= 0.
\end{aligned}
\right.
\tag{5}
\]
Diese Gleichungen stellen also eine irreducible Curve dar, deren Geschlecht $p$ durch die Gleichung (3) angegeben wird, und welche eine Gruppe von $n!$ eindeutigen Transformationen in sich besitzt, die zur Gruppe der Vertauschungen von $n$ Dingen holoedrisch isomorph ist. Das letztere wird durch die Form der Gleichungen (5) bestätigt: die eindeutigen Transformationen der Curve sind keine anderen, wie die durch die Vertauschungen der $y_1, y_2, \dots y_n$ definirten Collineationen des Raumes, in welchem unsere Curve liegt.

Das Beispiel der Curve (5) habe ich schon in dem letzten Paragraphen meiner Note über diejenigen algebraischen Gebilde, welche eindeutige Transformationen in sich zulassen, erwähnt. Ebenda habe ich als „Geschlecht“ einer Gruppe das niedrigste Geschlecht bezeichnet,

\origpage{427}
bei welchem algebraische Gebilde mit einem zu der Gruppe holoedrisch isomorphen System von eindeutigen Transformationen existiren. Ich bemerke, dass die Betrachtungen der vorigen Nummern in jedem Falle das Geschlecht einer gegebenen Gruppe zu bestimmen gestatten. In der That, seien
\[
S_1, \ S_2, \ \dots \ S_r
\tag{6}
\]
die Operationen der gegebenen Gruppe, $T_1$, $T_2$, $\dots T_{w-1}$ ein Theil derselben, welche zur Erzeugung der ganzen Gruppe ausreichen, endlich $T_w$ die Operation, welche der Gleichung
\[
T_1 T_2 \dots T_{w-1} T_w = 1
\]
genügt. Auf Grund der Operationen $T_1$, $T_2$, $\dots T_w$ kann man nun über der complexen Zahlenebene eine $r$-blättrige Riemann'sche Fläche construiren, welche eine zu (6) holoedrisch isomorphe Gruppe von eindeutigen Transformationen in sich besitzt. Sei $P$ das Geschlecht dieser Fläche. Um nun das Geschlecht der Gruppe (6) wirklich zu bestimmen, hat man alle Lösungen der diophantischen Gleichung
\[
\frac{2p - 2}{r} = \Bigl(1 - \frac{1}{k_1}\Bigr) + \Bigl(1 - \frac{1}{k_2}\Bigr) + \dots + (2\pi - 2),
\tag{7}
\]
welche den Bedingungen $p < P$, $k_1 \geqq 2$, $k_2 \geqq 2$, $\dots$, $\pi \geqq 0$ genügen, aufzusuchen. Die Anzahl dieser Lösungen ist offenbar eine endliche. Jede einzelne Lösung ist weiter daraufhin zu untersuchen, ob ihr eine Fläche $F$ entspricht, welche auf Grund der Operationen (6) nach den Vorschriften der Nr.~6 construirt ist. Wenn dies für keine Lösung der Fall ist, so ist $P$ das Geschlecht der Gruppe. Im andern Falle giebt diejenige der Flächen $F$, welche das kleinste Geschlecht besitzt, zugleich das Geschlecht der Gruppe. Discutirt man in dieser Weise z.~B. die Gruppe aller Vertauschungen bei fünf Dingen, so ergiebt sich für das Geschlecht dieser Gruppe der Werth $p = 4$.

\section*{III. Abschnitt.}

\subsection*{9.}

Besitzt eine Riemann'sche Fläche $F$, deren Geschlecht $p > 0$ ist, eine eindeutige Transformation $S$ in sich, so entspricht derselben immer eine lineare Transformation der $p$ überall endlichen Integrale $u_1, u_2, \dots u_p$. Entspricht nämlich vermöge $S$ der Stelle $P$ die Stelle $P'$ und bedeuten $u_1', u_2', \dots u_p'$ die Werthe der überall endlichen Integrale an der Stelle $P'$, so können wir diese Werthe auch als Functionen der Stelle $P$ ansehen. Als solche sind sie überall endliche Functionen, welche sich bei einem geschlossenen Weg, den $P$ durchläuft, um additive Constanten vermehren. Es stellen also $u_1', u_2', \dots u_p'$ als Functionen

\origpage{428}
der Stelle $P$ überall endliche Integrale vor. Es bestehen daher Gleichungen der Gestalt
\[
u_i' = \alpha_{i1} u_1 + \alpha_{i2} u_2 + \dots + \alpha_{ip} u_p + \alpha_i, \quad (i = 1, 2, \dots p),
\tag{1}
\]
wo $u_1, u_2, \dots u_p$ die Werthe der Integrale an der Stelle $P$ und die Coefficienten $\alpha$ von der Stelle $P$ unabhängige Grössen bedeuten. Es ist bequem, die Gleichungen (1) in differentieller Form zu schreiben:
\[
du_i' = \alpha_{i1} du_1 + \alpha_{i2} du_2 + \dots + \alpha_{ip} du_p, \quad (i = 1, 2, \dots p)
\tag{2}
\]
und wir sagen dieser Schreibweise entsprechend:
„\emph{Jeder Transformation $S$ der Fläche $F$ in sich entspricht eine homogene lineare Transformation} (2) \emph{der $p$ Differentiale erster Gattung.}“

Eine Gruppe von $r$ eindeutigen Transformationen der Fläche $F$ in sich entspricht daher eine Gruppe von homogenen linearen Transformationen der $p$ Differentiale erster Gattung. Und zwar sind beide Gruppen, sobald $p > 1$ ist, holoedrisch isomorph. Um dies zu beweisen, genügt es zu zeigen, dass die Transformation
\[
du_i' = du_i \quad (i = 1, 2, \dots p)
\tag{3}
\]
nur der identischen Transformation von $F$ entsprechen kann. Ist nun $S$ eine Transformation von $F$ in sich, welcher die Transformation (3) der Differentiale 1.~Gattung entspricht, so wird jedes Differential 1.~Gattung, welches in $P$ von der zweiten Ordnung verschwindet, zufolge (3) auch in $P'$ von der zweiten Ordnung Null. Daher muss entweder $P$ stets mit $P'$ zusammenfallen, also $S$ die identische Transformation sein, oder es muss die Fläche hyperelliptisch sein und $S$ in der Zuordnung je zweier Punkte $P, P'$ bestehen, in welchen die zweiwerthige Function der Fläche denselben Werth annimmt. Der letztere Fall ist aber ausgeschlossen, da bei dieser Zuordnung $du_i' = - du_i$ ist, während doch $du_i' = du_i$ sein soll.

Bei der Untersuchung einer Gruppe linearer Transformationen ist es eine erste fundamentale Aufgabe, die Wurzeln der charakteristischen Gleichung der einzelnen Transformation zu bestimmen. Wir wollen diese Aufgabe in dem hier vorliegenden Falle zu lösen versuchen. Es handelt sich also um folgende Frage:
„Gegeben ist eine (von der Identität verschiedene) eindeutige Transformation $S$ der Riemann'schen Fläche $F$ ($p > 1$) in sich. Welches sind die Wurzeln der charakteristischen Gleichung für die entsprechende homogene lineare Transformation (2) der Differentiale 1.~Gattung?“ Diese Frage soll nun zunächst in eine andere für die Beantwortung bequemere Form gebracht werden. Bekanntlich lässt sich eine lineare homogene Transformation von endlicher Ordnung oder Periode $n$, indem man die ursprünglichen Variabeln durch geeignete lineare Verbindungen derselben ersetzt, immer auf eine solche Form

\origpage{429}
bringen, dass die Transformation in der Multiplication der Variabeln mit $n^{\text{ten}}$ Einheitswurzeln besteht. Diese Einheitswurzeln sind dann die Wurzeln der charakteristischen Gleichung der Transformation. Wenn also $S$ die Periode $n$ besitzt, so können wir die $p$ Integrale $u_1, u_2, \dots u_p$ so wählen, dass (2) die Gestalt
\[
du_i' = \varepsilon_i\,du_i \quad (i = 1, 2, \dots p)
\tag{4}
\]
annimmt, wo $\varepsilon_1, \varepsilon_2, \dots \varepsilon_p$ $n^{\text{te}}$ Einheitswurzeln bedeuten. Unsere Frage lässt sich hiernach auch so formuliren:

\emph{Gegeben ist eine eindeutige Transformation $S$ der Riemann'schen Fläche $F(p > 1)$ in sich, welche jedem Punkte $P$ der Fläche $F$ einen bestimmten anderen Punkt $P'$ zuordnet und die Periode $n$ besitzt. Gegeben ist ferner eine $n^{\text{te}}$ Einheitswurzel $\varepsilon$. Wie viele linear unabhängige Integrale 1.~Gattung $u$ giebt es dann auf der Fläche $F$, welche der Bedingung}
\[
du' = \varepsilon\,du
\]
\emph{genügen, wo $u, u'$ die Werthe des Integrales in den Punkten $P, P'$ bedeuten?}

Die Zahl dieser linear unabhängigen Integrale giebt offenbar an, wie oft die Einheitswurzel $\varepsilon$ als Wurzel der charakteristischen Gleichung der Transformation (2) auftritt.

\subsection*{10.}

Wir betrachten die aus $S$ entspringende Gruppe
\[
S^0 = S_1, \quad S^1 = S_2, \quad S^2 = S_3, \quad \dots \quad S^{n-1} = S_n
\tag{1}
\]
von eindeutigen Transformationen der Fläche $F$ in sich. Vermöge der Transformationen (1) ordnen sich die Punkte von $F$ in Gruppen zu je $n$
\[
P_1, \ P_2, \ \dots \ P_n
\tag{2}
\]
an derart, dass vermöge $S$ jeder Punkt einer solchen Gruppe dem vorhergehenden, der erste Punkt dem letzten entspricht. Es finden nun hier ohne Weiteres die Betrachtungen der Nr.~6 Platz. Sei $\Phi$ die nach den Angaben jener Nummer hergestellte Fläche, welche $1-n$-deutig auf $F$ bezogen ist und zwar so, dass jedem Punkte von $\Phi$ die $n$ Punkte einer bestimmten Gruppe (2) auf $F$ entspricht. Seien ferner $a_1, a_2, \dots a_w$ die Punkte auf $\Phi$, deren zugehörige Gruppe (2) aus weniger als $n$ (also aus $\frac{n}{k}$ je $k$-fach zu rechnenden) Punkten besteht, dann können wir die Fläche $F$ $n$-blättrig über $\Phi$ ausbreiten und unter Beibehaltung der früheren Abkürzungen mit
\[
F = \begin{pmatrix}
l_1, & l_2, & \dots & l_w, & A_1, & B_1, & \dots & A_\pi, & B_\pi \\
T_1, & T_2, & \dots & T_w, & U_1, & V_1, & \dots & U_\pi, & V_\pi
\end{pmatrix}
\tag{3}
\]

\origpage{430}
bezeichnen. Hier sind die $T, U, V$ sämmtlich Potenzen von $S$, etwa
\[
\begin{gathered}
T_1 = S^{\mu_1}, \ \dots \ T_w = S^{\mu_w}, \quad U_1 = S^{\nu_1}, \quad V_1 = S^{\varrho_1}, \ \dots \ U_\pi = S^{\nu_\pi}, \\
V_\pi = S^{\varrho_\pi}.
\end{gathered}
\tag{4}
\]
Die Relation (7) in Nr.~6 geht über in
\[
T_1 T_2 \dots T_w = 1,
\]
oder, was dasselbe besagt,
\[
\mu_1 + \mu_2 + \dots + \mu_w \equiv 0 \pmod n.
\tag{5}
\]
Untersuchen wir jetzt die zu bestimmenden überall endlichen Integrale $u$, indem wir sie als Function der Stelle auf der Fläche $\Phi$ auffassen! Offenbar ist $u$ auf der Fläche $\Phi''$ (welche durch die Schnitte $l, A, B$ aus $\Phi$ entsteht) eindeutig und es erhält $du$ beim Uebertritt von der negativen zur positiven Seite der Schnitte
\[
l_1, \ l_2, \ \dots \ l_w, \ A_1, \ B_1, \ \dots \ A_\pi, \ B_\pi
\]
bezüglich die Factoren
\[
\varepsilon^{\mu_1}, \ \varepsilon^{\mu_2}, \ \dots \ \varepsilon^{\mu_w}, \ \varepsilon^{\nu_1}, \ \varepsilon^{\varrho_1}, \ \dots \ \varepsilon^{\nu_\pi}, \ \varepsilon^{\varrho_\pi}.{}^{*)}
\]
Da diese Eigenschaften für $u$ vollständig charakteristisch sind, so erkennen wir, dass sich unsere Frage auf die folgende reducirt:

\emph{Gegeben ist eine Riemann'sche Fläche $\Phi$, welche durch die Schnitte $l, A, B$ in eine einfach zusammenhängende Fläche $\Phi''$ zerschnitten ist. Wie viele linear unabhängige überall endliche Functionen $u$ giebt es, welche auf $\Phi''$ eindeutig sind und die Eigenschaft besitzen, dass das Differential $du$ beim Uebertritt von der negativen auf die positive Seite der Schnitte}
\[
l_1, \ l_2, \ \dots \ l_w, \ A_1, \ B_1, \ \dots \ A_\pi, \ B_\pi
\]
\emph{bezüglich die gegebenen Factoren}
\[
\gamma_1, \ \gamma_2, \ \dots \ \gamma_w, \ \alpha_1, \ \beta_1, \ \dots \ \alpha_\pi, \ \beta_\pi
\]
\emph{erhalten?}

Die gegebenen Factoren sind $n^{\text{te}}$ Einheitswurzeln, zwischen welchen (nach 5) die Beziehung
\[
\gamma_1 \gamma_2 \dots \gamma_w = 1
\tag{6}
\]
besteht.

\textbf{*)} D.~h. Sind $u^+$ und $u^-$ die Werthe, welche die auf $\Phi''$ eindeutig ausgebreitete Function $u$ in gegenüberliegenden Punkten des Schnittes $l_1$ besitzt, so ist längs des ganzen Schnittes $l_1$
\[
u^+ = \varepsilon^{\mu_1} \cdot u^- + \text{const.}
\]
Entsprechend für die anderen Schnitte.

\origpage{431}
\subsection*{11.}

Unsere Frage führt also auf die Bestimmung solcher Functionen des Ortes einer Riemann'schen Fläche $\Phi$, welche sich bei Durchlaufung eines geschlossenen Weges bis auf eine multiplicative Constante reproduciren. Die Theorie dieser Functionen und ihrer Integrale bildet eine Verallgemeinerung der Theorie der algebraischen Functionen und ihrer Integrale, welch' letztere dem Fall entspricht, in welchem die multiplicativen Constanten sämmtlich gleich 1 sind. Eine eingehende Untersuchung der in Rede stehenden Functionen gedenke ich demnächst zu veröffentlichen. Hier beschränke ich mich darauf, diejenigen Sätze anzugeben, welche für den vorliegenden Zweck erforderlich sind. Zuvor bemerke ich noch, dass der specielle Fall, in welchem die zu untersuchenden Functionen auf der Fläche $\Phi$ unverzweigt sind, schon mehrfach behandelt ist. Bereits Riemann hat solche Functionen in Nr.~26 seiner „Theorie der Abel'schen Functionen“ betrachtet.*) (Die dort eingeführte Beschränkung, dass die multiplicativen Constanten Einheitswurzeln seien, ist unwesentlich.) Sodann hat Herr Prym im 70.~Bande von Crelle's Journal einige Sätze über die Bestimmung dieser Functionen aus ihren Periodicitätseigenschaften ohne Beweis mitgetheilt. Endlich ist die Preisarbeit „Sur les intégrales de fonctions à multiplicateurs et leur application au développement des fonctions abéliennes en séries trigonométriques“ (Acta mathematica, Bd.~13) von Herrn Appell zu erwähnen, deren erster und zweiter Theil sich auf dieselben Functionen beziehen.**) Ich will nun zunächst die zu betrachtenden Functionen genau definiren. Es sei $\Phi$ eine Riemann'sche Fläche vom Geschlecht $\pi$, ferner
\[
a_1, \ a_2, \ \dots \ a_w
\]
willkürlich auf $\Phi$ angenommene Punkte. Wir zerschneiden $\Phi$ durch die Schnitte
\[
l_1, \ l_2, \ \dots \ l_w, \ A_1, \ B_1, \ \dots \ A_\pi, \ B_\pi
\]
in die einfach zusammenhängende Fläche $\Phi''$. Die Functionen $f$ des Ortes auf $\Phi$, welche wir betrachten wollen, sind dann durch folgende Eigenschaften charakterisirt:
1) Verzweigungspunkte von $f$ finden sich nur unter den Punkten $a_1, a_2, \dots a_w$, so dass die einzelnen Zweige von $f$ auf der zerschnittenen Fläche $\Phi''$ eindeutig sind.

\textbf{*)} Man sehe auch die interessante (klein gedruckte) Bemerkung auf pag.~364 der gesammelten Werke von Riemann.

\textbf{**)} Die Untersuchungen Appell's lassen sich durch Einführung des oben gewählten Schnittsystems vereinfachen. Es kommen dann nämlich alle auf die Schnitte $c$ bezüglichen Betrachtungen in Wegfall.

\origpage{432}
2) Die Function $f$ besitzt, abgesehen von den Punkten $a_1, a_2, \dots a_w$, keine anderen singulären Punkte als Pole.

3) In der Umgebung des Punktes $a_i$ bleibt
\[
t^{-\delta_i} \cdot f
\]
endlich, eindeutig und von Null verschieden, wenn $t$ eine im Punkte $a_i$ von der ersten Ordnung verschwindende Function und $\delta_i$ eine passend gewählte Constante bedeutet.

4) Jeder auf $\Phi''$ eindeutig ausgebreitete Zweig der Function $f$ erhält beim Uebertritt von der negativen zur positiven Seite eines der Schnitte
\[
l_1, \ l_2, \ \dots \ l_w, \ A_1, \ B_1, \ \dots \ A_\pi, \ B_\pi
\]
bezüglich die gegebenen Factoren
\[
\gamma_1, \ \gamma_2, \ \dots \ \gamma_w, \ \alpha_1, \ \beta_1, \ \dots \ \alpha_\pi, \ \beta_\pi.
\]
Das System aller Functionen $f$, welches diesen Bedingungen genügt, bezeichnen wir in Riemann'scher Weise mit
\[
f = \begin{pmatrix}
l_1, & l_2, & \dots & l_w, & A_1, & B_1, & \dots & A_\pi, & B_\pi \\
\gamma_1, & \gamma_2, & \dots & \gamma_w, & \alpha_1, & \beta_1, & \dots & \alpha_\pi, & \beta_\pi
\end{pmatrix}. \tag{1}
\]
Da in der Umgebung des Punktes $a_i$ die Entwicklung von $f$ die Gestalt
\[
f = t^{\delta_i}(c_0 + c_1 t + c_2 t^2 + \dots) \tag{2}
\]
besitzen soll, so ist klar, dass
\[
\gamma_1 = e^{2\pi i\delta_1}, \quad \gamma_2 = e^{2\pi i\delta_2}, \quad \dots \quad \gamma_w = e^{2\pi i\delta_w} \tag{3}
\]
sein muss. Betrachten wir nun das in positivem Sinne durch die Begrenzung von $\Phi''$ erstreckte Integral
\[
\frac{1}{2\pi i}\int d\lg f,
\]
so ist dasselbe gleich $N - U$, wo $N$ die Zahl der Null-, $U$ die Zahl der Unendlichkeitsstellen von $f$ bedeutet. Andererseits reducirt sich jenes Integral auf diejenigen Theilintegrale, welche um die Punkte $a_1$, $a_2$, \dots $a_w$ in negativem Sinne herumführen. Die letzteren Theilintegrale besitzen nach (2) die Werthe $-\delta_1$, $-\delta_2$, \dots, $-\delta_w$. Es ist also
\[
U = N + \delta_1 + \delta_2 + \dots + \delta_w. \tag{4}
\]
Hieraus geht hervor, dass $\delta_1 + \delta_2 + \dots + \delta_w$ nothwendig eine ganze Zahl und folglich
\[
\gamma_1 \gamma_2 \dots \gamma_w = 1 \tag{5}
\]
ist.*) Die Factoren $\gamma, \alpha, \beta$ dürfen also nicht vollkommen willkürlich

\textbf{*)} Dasselbe folgt auch durch Betrachtung eines Umlaufs um den Punkt $O$, von welchem die Schnitte $l$, $A$, $B$ ausgehen.

\origpage{433}
gewählt werden, sondern müssen der Bedingung (5) genügen. Ausserdem darf natürlich keiner der Factoren gleich Null sein. Es findet aber keine weitere Bedingung für diese Factoren statt. In der That besteht der Satz:

\emph{Zu jedem Factorensystem}
\[
\gamma_1, \ \gamma_2, \ \dots \ \gamma_w, \ \alpha_1, \ \beta_1, \ \dots \ \alpha_\pi, \ \beta_\pi,
\]
\emph{welches den Bedingungen $\gamma_1 \gamma_2 \dots \gamma_w = 1$ und $\alpha_1 \beta_1 \dots \alpha_\pi \beta_\pi \neq 0$ genügt, übrigens aber ganz beliebig gewählt ist, gehört stets ein System von Functionen}
\[
f = \begin{pmatrix}
l_1, & l_2, & \dots & l_w, & A_1, & B_1, & \dots & A_\pi, & B_\pi \\
\gamma_1, & \gamma_2, & \dots & \gamma_w, & \alpha_1, & \beta_1, & \dots & \alpha_\pi, & \beta_\pi
\end{pmatrix}.
\]

Man kann nämlich stets auf mannigfaltige Weise eine Summe $\varSigma$ von Abel'schen Integralen dritter Gattung bilden, so dass
\[
f = e^{\varSigma} \tag{6}
\]
eine alle erforderlichen Eigenschaften besitzende Function darstellt. Man erkennt auch leicht, dass durch den Ausdruck (6) die allgemeinste Function des Systems dargestellt wird.*)

\subsection*{12.}

Zu jedem Functionensystem
\[
\begin{pmatrix}
l_1, & l_2, & \dots & l_w, & A_1, & B_1, & \dots & A_\pi, & B_\pi \\
\gamma_1, & \gamma_2, & \dots & \gamma_w, & \alpha_1, & \beta_1, & \dots & \alpha_\pi, & \beta_\pi
\end{pmatrix} \tag{1}
\]
gehört ein System von Integralfunktionen, welches durch die Integrale
\[
\int f\,dz
\]
gebildet wird. Dabei bedeutet $f$ irgend eine Function des Systems (1) und $z$ eine eindeutige algebraische Function oder ein Abel'sches Integral der Fläche $\Phi$.

Ich will hier nur die überall endlichen unter diesen Integralfunctionen näher betrachten. Es handelt sich vor Allem darum (vgl. Nr.~10), die Zahl der linear unabhängigen dieser Functionen zu bestimmen. Hierbei werde ich die Multiplicatoren je in folgende Form setzen:

\textbf{*)} Ich bemerke hier ausdrücklich, dass das Geschlecht $\pi$ der Fläche $\Phi$ beliebig angenommen wird. Es ist also der Fall $\pi = 0$ nicht ausgeschlossen. In diesem Falle kommen nur die Schnitte $A$, $B$ in Fortfall. Die Zahl $w$ darf man stets positiv voraussetzen, da man immer beliebig viele Schnitte $l$ aufnehmen kann, denen man den Factor $\gamma = 1$ zuordnet, ohne dadurch das Functionssystem $f$ zu ändern.

\origpage{434}
\[
\gamma_1 = e^{-2\pi i\varkappa_1}, \quad \gamma_2 = e^{-2\pi i\varkappa_2}, \quad \dots \quad \gamma_w = e^{-2\pi i\varkappa_w}, \tag{2}
\]
wo die Zahlen $\varkappa_1, \varkappa_2, \dots \varkappa_w$ durch die Festsetzung unzweideutig bestimmt sind, dass ihre reellen Bestandtheile zwischen 0 und 1 liegen sollen, die obere Grenze 1 ausgeschlossen. Es sei nun $v$ eine der überallendlichen Integralfunctionen, ferner $P$ eine Stelle der Fläche $\Phi$ und $t$ eine in $P$ von der ersten Ordnung verschwindende Function. Gehört nun $P$ nicht zu den Stellen $a_1, a_2, \dots a_w$, so lautet die Entwicklung von $v$ an der Stelle $P$
\[
v = c + c' t^r + c'' t^{r+1} + \dots \quad (c' \neq 0). \tag{3}
\]
Wir sagen dann (genau wie für die eindeutigen algebraischen Functionen und ihre Integrale), dass die Stelle $P$ eine $(r-1)$-fach zu rechnende Nullstelle von $dv$ sei.

Fällt aber $P$ mit $a_i$ zusammen, so haben wir eine Entwicklung der Gestalt\ednote{In the following display the typesetter omitted the factor $t$ after $c''$; the series in powers of $t$ should read $c' + c'' t + \dots$. Reproduced here as printed.}
\[
v = c + t^{-\varkappa_i + r}(c' + c'' + \dots) \quad (c' \neq 0) \tag{4}
\]
und wir wollen dann $a_i$ als eine $(r-1)$-fache Nullstelle von $dv$ rechnen. Die Zahlen $r$ in (3) und (4) sind positive ganze Zahlen, weil $v$ überall endlich sein soll. Dies festgesetzt, beweisen wir zunächst den Satz:

„\emph{Die Gesammtzahl der Nullstellen des Differentials $dv$ beträgt}
\[
2\pi - 2 + \varkappa_1 + \varkappa_2 + \dots + \varkappa_w\text{.“}
\]
In der That, sei $z$ eine eindeutige algebraische Function auf $\Phi$, welche $m$ Mal unendlich wird. Dann ist $\frac{dv}{dz}$ eine Function $f$, auf die wir die Gleichung (4) von Nr.~1\ednote{Printed Nr. 1 instead of Nr. 11; equation (4) was established in §11.} anwenden können. Nun wird $\frac{dv}{dz}$ unendlich nur an den Nullstellen von $dz$, deren es bekanntlich $2m+2\pi-2$ giebt. Ferner wird $\frac{dv}{dz}$ an den $m$ Unendlichkeitsstellen von $z$ je von der zweiten Ordnung Null und an jeder von den Stellen $a_1, a_2, \dots a_w$ verschiedenen Nullstelle von $dv$ so oft Null, als dieselbe nach unserer Festsetzung zu rechnen ist.

An der Stelle $a_i$ lautet die Entwicklung von $\frac{dv}{dz}$
\[
\frac{dv}{dz} = t^{-\varkappa_i + (r-1)} (c_0' + \dots),
\]
wenn wir annehmen, was offenbar gestattet ist, dass die Stellen $a_i$ nicht zu den Nullstellen von $dz$ und ebenso wenig zu den Unendlichkeitsstellen von $z$ gehören.

Daher giebt die Gleichung (4) von Nr.~11
\[
2m + 2\pi - 2 = 2m + \sum(r-1) - \varkappa_1 - \varkappa_2 - \dots - \varkappa_w,
\]
woraus die Richtigkeit unseres Satzes hervorgeht. Wir ziehen aus

\origpage{435}
diesem Satze noch eine Folgerung. Da die Anzahl der Nullstellen von $dv$ nicht negativ werden kann, so zeigt unser Satz, dass in den Fällen
\[
\pi = 0, \quad \varkappa_1 + \varkappa_2 + \dots + \varkappa_w = 0
\]
und
\[
\pi = 0, \quad \varkappa_1 + \varkappa_2 + \dots + \varkappa_w = 1
\]
keine überall endliche Integralfunction $v$ existiren kann.

Wir wollen nunmehr zeigen, wie man alle Integralfunctionen $v$ bestimmen kann unter der Voraussetzung, dass man eine dieser Functionen kennt.

Es sei $v_0$ diese eine Integralfunction, welche wir als bekannt voraussetzen. Ist dann $v$ eine beliebige der Integralfunctionen, so wird offenbar
\[
\frac{dv}{dv_0} = Z \tag{5}
\]
eine eindeutige algebraische Function der Fläche sein, welche nur an den $2\pi - 2 + \varkappa_1 + \varkappa_2 + \dots + \varkappa_w$ Nullstellen von $dv_0$ unendlich werden kann. Ist umgekehrt $Z$ eine derartige Function, so wird
\[
v = \int Z\,dv_0 \tag{6}
\]
eine Integralfunction sein. Nun giebt es aber nach dem Riemann-Roch'schen Satze genau
\[
2\pi - 2 + \varkappa_1 + \varkappa_2 + \dots + \varkappa_w - \pi + 1 = \varkappa_1 + \varkappa_2 + \dots + \varkappa_w + \pi - 1
\]
linear unabhängige Functionen $Z$, wenn wir von einem sogleich näher zu betrachtenden Ausnahmefall absehen. Es muss also ebenso viele linear unabhängige Integralfunctionen $v$ geben, welche sämmtlich nach Formel (6) aus der einen $v_0$ abgeleitet werden können.

Der erwähnte Ausnahmefall tritt ein, wenn die Nullstellen von $dv_0$ zugleich die Nullstellen eines Differentials erster Gattung $du$ sind. Dies kann nur dann stattfinden, wenn
\[
\varkappa_1 + \varkappa_2 + \dots + \varkappa_w = 0,
\]
also auch die reellen Bestandtheile von $\varkappa_1, \varkappa_2, \dots \varkappa_w$ sämmtlich Null und daher die Factoren $\gamma_1, \gamma_2, \dots \gamma_w$ sämmtlich reell sind. Die Anzahl der linear unabhängigen Differentiale erster Gattung, welche in den Nullstellen von $dv_0$ verschwinden, kann nur gleich 1 sein, weil die Anzahl dieser Stellen $2\pi - 2$ beträgt. Die oben angegebene Zahl ist dann also zu ersetzen durch
\[
\varkappa_1 + \varkappa_2 + \dots + \varkappa_w + \pi = \pi.
\]
Um diesen Ausnahmefall noch näher zu charakterisiren, betrachten wir die Function $\frac{dv_0}{du}$. Dieselbe gehört zu dem Functionssysteme $f$; sie wird nirgends Null oder unendlich, ausgenommen etwa an den Stellen $a_1, \dots a_w$, wo ihr Verhalten durch die Entwicklung

\origpage{436}
\[
\frac{dv_0}{du} = t^{-\varkappa_i}(c_1 + \dots) \quad (\text{an der Stelle } a_i),
\]
angezeigt wird. Wir schliessen hieraus, dass $\log\frac{dv_0}{du}$ ein Integral 3.~Gattung ist, welches an der Stelle $a_i$ wie $-\varkappa_i\log t$ unstetig wird. Bezeichnen wir daher mit $\Pi_{a, a_i}$ ein Integral 3.~Gattung, welches an der willkürlich gewählten Stelle $a$ wie $\log t$, an der Stelle $a_i$ wie $-\log t$ unstetig wird, so ist
\[
\log\frac{dv_0}{du} - \varkappa_1 \Pi_{a, a_1} - \varkappa_2 \Pi_{a, a_2} - \dots - \varkappa_w \Pi_{a, a_w} = u_1
\]
ein Integral erster Gattung (wenn wir eine Constante unter den Begriff eines Integrals 1.~Gattung einschliessen). Daher kommt
\[
\frac{dv_0}{du} = \mathrm{E}, \tag{7}
\]
unter $\mathrm{E}$ den Ausdruck
\[
\mathrm{E} = e^{\varkappa_1 \Pi_{a, a_1} + \varkappa_2 \Pi_{a, a_2} + \dots + \varkappa_w \Pi_{a, a_w} + u_1} \tag{8}
\]
verstanden. Da umgekehrt die $\pi$ in der Form $\displaystyle\int \mathrm{E}\,du$ enthaltenen linear unabhängigen Integralfunctionen überall endlich sind, so können wir folgenden Satz aussprechen:

„\emph{Der Ausnahmefall tritt stets und nur dann ein, wenn das Factorensystem $\gamma, \alpha, \beta$ übereinstimmt mit den Factoren, welche ein Ausdruck $\mathrm{E}$ der Gestalt} (8) \emph{besitzt, in welchem $\varkappa_1, \varkappa_2, \dots \varkappa_w$ rein imaginäre Grössen bedeuten.“}

\subsection*{13.}

Es erübrigt noch zu untersuchen, in welchen Fällen mindestens eine überall endliche Integralfunction $v$ existirt. Sei $f$ irgend eine Function des Systems (1) (Nr.~12), $z$ eine eindeutige algebraische Function der Fläche $\Phi$. Unter der Annahme, dass $v$ eine überall endliche Integralfunction sei, betrachten wir den Quotienten
\[
\frac{dv}{f\,dz} = Z. \tag{1}
\]
Derselbe stellt eine eindeutige algebraische Function der Fläche $\Phi$ vor, welche, abgesehen von den Stellen $a_1, a_2, \dots a_w$, nur unstetig werden kann

1) an den $N$ Nullstellen von $f$,

2) an den $2m+2\pi-2$ Nullstellen von $dz$.

An der Stelle $a_i$ ist
\[
f = t^{\delta_i}(c_0 + c_1 t + \dots),
\]
und wir dürfen voraussetzen, dass die ganzen Zahlen $\delta_i + \varkappa_i$ positiv

\origpage{437}
sind, da wir andernfalls $f$ mit einer an den Stellen $a_1, a_2, \dots a_w$ verschwindenden eindeutigen algebraischen Function $z_1$ multipliciren, also $z_1 \cdot f$ an Stelle von $f$ setzen können.

Es wird dann $Z$

3) an der Stelle $a_i$ unendlich höchstens von der Ordnung $\delta_i + \varkappa_i$.

Ausserdem verschwindet $Z$

4) an den $m$ Unendlichkeitsstellen von $z$ je von der zweiten Ordnung, und an den $U$ Unendlichkeitsstellen von $f$.
Umgekehrt: Ist $Z$ eine eindeutige algebraische Function, welche den Bedingungen 1), 2), 3), 4) genügt, so wird
\[
v = \int Z f\,dz
\]
eine überall endliche Integralfunction $v$ darstellen.

Nun enthalten die Functionen $Z$, welche den Bedingungen 1), 2), 3) genügen, nach dem Riemann-Roch'schen Satze
\[
N + 2m + 2\pi - 2 + (\delta_1 + \varkappa_1) + (\delta_2 + \varkappa_2) + \dots + (\delta_w + \varkappa_w) - \pi + 1
\]
willkürliche Constante linear und homogen. Für diese Constante liefern die Bedingungen 4) $2m + U$ lineare homogene Bedingungsgleichungen.

Daher enthält die allgemeinste Function $Z$, welche allen Bedingungen 1)---4) genügt, mindestens noch
\[
N + 2m + 2\pi - 2 + (\delta_1 + \varkappa_1) + \dots + (\delta_w + \varkappa_w) - \pi + 1 - (2m + U)
\]
willkürliche Constante linear und homogen. Diese Zahl reducirt sich aber nach Gleichung (4) in Nr.~11 auf
\[
\varkappa_1 + \varkappa_2 + \dots + \varkappa_w + \pi - 1 \tag{2}
\]
und ebenso viele überall endliche Integralfunctionen $v$ giebt es also \emph{mindestens}. Sobald nun die Zahl (2) grösser als Null ist, treten die Betrachtungen der vorigen Nummer in Kraft. Die einzigen Fälle, in welchen jene Zahl nicht grösser als Null ist, sind
\[
\begin{aligned}
\pi &= 1, & \varkappa_1 + \varkappa_2 + \dots + \varkappa_w &= 0, \\
\pi &= 0, & \varkappa_1 + \varkappa_2 + \dots + \varkappa_w &= 1, \\
\pi &= 0, & \varkappa_1 + \varkappa_2 + \dots + \varkappa_w &= 0.
\end{aligned}
\]
Der erste dieser Fälle subsumirt sich unter den pag.~435 besprochenen Ausnahmefall. In den beiden letzten Fällen giebt es keine überall endliche Integralfunction $v$, wie wir schon oben bemerkt haben. Indem wir dieses berücksichtigen, gelangen wir zu folgendem allgemein gültigen Satze, welcher in allen Fällen die Anzahl der überall endlichen Integralfunctionen, welche zu einem Functionensysteme\ednote{In the following matrix the entry above $\beta_1$ is printed $\alpha_j$ with subscript $j$ instead of $1$; an apparent compositor misprint for $\alpha_1$. Reproduced here as printed.}
\[
f = \begin{pmatrix}
l_1, & l_2, & \dots & l_w, & A_1, & B_1, & \dots & A_\pi, & B_\pi \\
\gamma_1, & \gamma_2, & \dots & \gamma_w, & \alpha_j, & \beta_1, & \dots & \alpha_\pi, & \beta_\pi
\end{pmatrix}
\]
gehören, bestimmen lehrt:

\origpage{438}
\emph{Man setze}
\[
\gamma_1 = e^{-2\pi i\varkappa_1}, \quad \gamma_2 = e^{-2\pi i\varkappa_2}, \quad \dots \quad \gamma_w = e^{-2\pi i\varkappa_w}, \tag{3}
\]
\emph{wo die Grössen $\varkappa_1, \varkappa_2, \dots \varkappa_w$ durch die Forderung eindeutig bestimmt sind, dass ihre reellen Bestandtheile kleiner als 1 und nicht negativ sein sollen*).}

\emph{Dann beträgt die Anzahl der linear unabhängigen überall endlichen Integralfunctionen $v$}
\[
\varkappa_1 + \varkappa_2 + \dots + \varkappa_w + \pi - 1. \tag{4}
\]
\emph{Ausgenommen ist nur der eine Fall, in welchem die Multiplicatoren $\gamma$ sämmtlich reell, also $\varkappa_1, \varkappa_2, \dots \varkappa_w$ rein imaginär und}
\[
\varkappa_1 + \varkappa_2 + \dots + \varkappa_w = 0
\]
\emph{ist, und zugleich die Factoren $\alpha, \beta, \gamma$ übereinstimmen mit den Factoren eines Ausdrucks der Gestalt}
\[
\mathrm{E} = e^{\varkappa_1 \Pi_{a, a_1} + \varkappa_2 \Pi_{a, a_2} + \dots + \varkappa_w \Pi_{a, a_w} + u}.
\]
\emph{In diesem Falle ist die Anzahl} (4) \emph{zu ersetzen durch die Anzahl}
\[
\pi. \tag{5}
\]

Ich bemerke noch, dass $\Pi_{a, a_1}$ ein Integral dritter Gattung bedeutet, welches bei $a$ unstetig wird wie $\log t$, bei $a_1$ unstetig wie $-\log t$. Wenn $\Phi$ das Geschlecht $\pi = 0$ besitzt und $z$ diejenige (bis auf eine lineare Substitution bestimmte) Function bezeichnet, welche jeden Werth nur an einer Stelle annimmt, wird also $\Pi_{a, a_1} = \log\frac{z - z_a}{z - z_{a_1}} + \text{const.}$ sein, wo $z_a$, $z_{a_1}$ die Werthe von $z$ an den Stellen $a$ und $a_1$ bedeuten. Die im Exponenten von $e$ auftretende Grösse $u$ bezeichnet ein überall endliches Integral, wobei indessen auch eine Constante als ein solches anzusehen ist. Im Falle $\pi = 0$ wird $u$ stets eine Constante sein.

Für den speciellen Fall, in welchem die Factoren $\gamma$ sämmtlich gleich 1 sind, in welchem es sich also um das Functionensystem
\[
f = \begin{pmatrix}
A_1, & B_1, & \dots & A_\pi, & B_\pi \\
\alpha_1, & \beta_1, & \dots & \alpha_\pi, & \beta_\pi
\end{pmatrix}
\]
handelt, ergiebt unser Satz ein von Herrn Appell (l.~c. pag.~25, 26) gefundenes Resultat.

\subsection*{14.}

Wir haben jetzt alle Mittel, um die in Nr.~10 aufgeworfene Frage zu beantworten. Es handelte sich dort um die Bestimmung der Anzahl

\textbf{*)} Es kann wohl kein Missverständniss herbeiführen, dass in den Formeln (3) $\pi$ die Ludolph'sche Zahl bedeutet, während andererseits in den Formeln (4) und (5) wie überall in diesen Betrachtungen $\pi$ das Geschlecht der Riemann'schen Fläche $\Phi$ bezeichnet.

\origpage{439}
der linear unabhängigen überall endlichen Integralfunctionen des Systemes
\[
f = \begin{pmatrix}
l_1, & l_2, & \dots & l_w, & A_1, & B_1, & \dots & A_\pi, & B_\pi \\
\varepsilon^{\mu_1}, & \varepsilon^{\mu_2}, & \dots & \varepsilon^{\mu_w}, & \varepsilon^{\nu_1}, & \varepsilon^{\varrho_1}, & \dots & \varepsilon^{\nu_\pi}, & \varepsilon^{\varrho_\pi}
\end{pmatrix},
\]
wobei $\varepsilon$ eine $n^{\text{te}}$ Einheitswurzel bezeichnet. Wir setzen
\[
\varepsilon = e^{-2\pi i \cdot \frac{k}{n}} \tag{1}
\]
wo $k$ eine Zahl der Reihe $0, 1, 2, \dots n-1$ bedeutet. Dann wird
\[
\gamma_1 = e^{-2i\pi\frac{k\mu_1}{n}}, \quad \gamma_2 = e^{-2i\pi\frac{k\mu_2}{n}}, \quad \dots \quad \gamma_w = e^{-2i\pi\frac{k\mu_w}{n}}
\]
und daher kommt,
\[
\varkappa_1 = \frac{[k\mu_1]}{n}, \quad \varkappa_2 = \frac{[k\mu_2]}{n}, \quad \dots \quad \varkappa_w = \frac{[k\mu_w]}{n},
\]
wenn allgemein mit $[k\mu_i]$ der kleinste nicht negative Rest von $k\mu_i\,(\text{mod. } n)$ bezeichnet wird. Die Anzahl der linear unabhängigen Integralfunctionen beträgt also
\[
A_k = \frac{[k\mu_1] + [k\mu_2] + \dots + [k\mu_w]}{n} + \pi - 1, \tag{2}
\]
wenn wir zunächst von dem Ausnahmefall absehen.

Der Ausnahmefall tritt nur ein, wenn

1) die reellen Bestandtheile von $\varkappa_1, \varkappa_2, \dots \varkappa_w$ verschwinden, also da $\varkappa_1, \varkappa_2, \dots \varkappa_w$ hier reelle Zahlen sind, wenn $\varkappa_1 = \varkappa_2 = \dots = \varkappa_w = 0$, folglich $\gamma_1 = \gamma_2 = \dots = \gamma_w = 1$ ist,

und wenn

2) die Factoren $\gamma, \alpha, \beta$ übereinstimmen mit den Factoren eines Ausdrucks $e^u$, wo $u$ ein überall endliches Integral der Fläche $\Phi$ bezeichnet. Da aber $\gamma, \alpha, \beta$ $n^{\text{te}}$ Einheitswurzeln sind, kann dies nur stattfinden, wenn $u$ eine Constante ist. Denn $e^{nu}$ ist eine überall endliche eindeutige algebraische Function der Fläche $\Phi$. Ist aber $e^u$ eine Constante, so sind die Factoren, welche $e^u$ beim Ueberschreiten eines Querschnitts annimmt, sämmtlich gleich 1. Folglich tritt der Ausnahmefall nur ein, wenn die Factoren $\gamma, \alpha, \beta$ sämmtlich gleich 1 sind, wenn es sich also um die überall endlichen Integrale der Fläche $\Phi$ handelt.
Für diese Integrale ist die Einheitswurzel $\varepsilon = 1$, folglich $k = 0$ und die Formel (2) erleidet also nur im Falle $k = 0$ eine Abänderung. Sie ist dann zu ersetzen durch
\[
A_0 = \pi. \tag{3}
\]

\origpage{440}
\subsection*{15.}

Wir kehren jetzt zur Betrachtung der Fläche $F$ zurück, bezeichnen wie in Nr.~10 mit $S$ die eindeutige Transformation von der Periode $n$, welche die Fläche $F$ zulässt, und mit
\[
P_1, \ P_2, \ \dots \ P_n \tag{1}
\]
die Gruppe von $n$ Punkten, welche aus einem willkürlich gewählten Punkt $P_1$ der Fläche $F$ vermöge der Transformation $S$ hervorgehen.

Die Punktgruppen (1) sind auf die Stellen der Fläche $\Phi$ eindeutig umkehrbar bezogen und die Fläche $F$ war dem entsprechend darstellbar in Gestalt einer $n$-blättrig über $\Phi$ ausgebreiteten Fläche:
\[
F = \begin{pmatrix}
l_1, & l_2, & \dots & l_w, & A_1, & B_1, & \dots & A_\pi, & B_\pi \\
T_1, & T_2, & \dots & T_w, & U_1, & V_1, & \dots & U_\pi, & V_\pi
\end{pmatrix}.
\]
Der Punkt $a_i$ der Fläche $\Phi$, nach welchem der Schnitt $l_i$ läuft, ist einer derjenigen Punkte, denen auf der Fläche $F$ eine Gruppe (1) von $\frac{n}{k_i}$ je $k_i$-fach zu zählenden Punkten entspricht. Der Substitution $T_i = S^{\mu_i}$ entsprechend, hängen nun die mit $S_1, S_2, \dots S_n$ bezeichneten Blätter der Fläche $F$ an der Stelle $a_i$ so zusammen, dass bei einem Umlauf um $a_i$ allgemein $S_k$ in $S_{k+\mu_i}$ übergeht, wo die Indices $(\text{mod. } n)$ zu nehmen sind. Bestimmt man hiernach die Zahl der Blätter, welche bei $a_i$ je im Cyklus zusammenhängen --- eine Zahl, die offenbar mit $k_i$ identisch ist --- so erhält man den Satz:

„\emph{Die Zahl $k_i$ ist die kleinste positive Lösung der Congruenz}
\[
k_i\mu_i \equiv 0 \pmod n\text{.“}
\]
Die Anzahl der linear unabhängigen überall endlichen Integrale $u$ der Fläche $F$, welche der Bedingung
\[
du' = \varepsilon\,du
\]
genügen, wo $u, u'$ die Werthe des Integrals an den Stellen $P_1, P_2$ und
\[
\varepsilon = e^{-2i\pi\frac{k}{n}}
\]
eine gegebene $n^{\text{te}}$ Einheitswurzel bedeuten, beträgt nun nach der vorigen Nummer
\[
A_k = \frac{[k\mu_1] + [k\mu_2] + \dots + [k\mu_w]}{n} + \pi - 1, \tag{2}
\]
wenn $k$ einen der Werthe $1, 2, \dots n-1$ besitzt, dagegen
\[
A_0 = \pi, \tag{3}
\]
wenn $k = 0$, also $\varepsilon = 1$ ist.

Die Gesammtzahl der Integrale
\[
A_0 + A_1 + A_2 + \dots + A_{n-1}
\]

\origpage{441}
muss gleich $p$ sein und wir wollen zum Schluss zeigen, dass diese Thatsache mit den für die Zahlen $A$ gefundenen Werthen (2) und (3) im Einklang steht.

Das Zeichen $[r]$, welches in (2) auftritt, bedeutete den kleinsten nicht negativen Rest der Zahl $r$ nach dem Modul $n$. Hieraus geht unmittelbar hervor, dass
\[
[r] = [r'], \quad \text{für } r \equiv r' \pmod n, \tag{4}
\]
\[
[r] = n - [r'], \quad \text{für } r \equiv -r' \pmod n, \text{ wenn zugleich } r \text{ nicht} \equiv 0 \pmod n. \tag{5}
\]
Ist $r \equiv 0 \pmod n$, so wird nach Definition $[r] = 0$.

Dies vorausgeschickt, betrachten wir die Summe
\[
\begin{aligned}
A_k + A_{n-k} = 2(\pi - 1) &+ \frac{[k\mu_1] + \dots + [k\mu_w]}{n} \\
&+ \frac{[(n-k)\mu_1] + \dots + [(n-k)\mu_w]}{n},
\end{aligned} \tag{6}
\]
wo $k$ eine der Zahlen $1, 2, \dots n-1$ bedeutet. Der Theil
\[
\frac{[k\mu_1] + [(n-k)\mu_1]}{n}
\]
dieser Summe ist nach (5) gleich 1, ausgenommen den Fall, wo $k\mu$ durch $n$ theilbar ist, in welchem jener Theil gleich Null ist. Nun ist aber
\[
k\mu_1 \equiv 0 \pmod n
\]
dann und nur dann, wenn $k$ durch $k_1$ theilbar ist, denn nach dem oben Bemerkten ist $k_1$ die kleinste positive Lösung der Congruenz $k_1\mu_1 \equiv 0 \pmod n$.

Indem man die entsprechende Ueberlegung für $\frac{[k\mu_i] + [(n-k)\mu_i]}{n}$ anstellt, erkennt man, dass
\[
A_k + A_{n-k} = 2(\pi - 1) + w - D_k \tag{7}
\]
ist, wo $D_k$ angiebt, wie viele der Zahlen
\[
k_1, \ k_2, \ \dots \ k_w
\]
in die Zahl $k$ aufgehen.*)

\textbf{*)} Ist $n$ eine Primzahl, so sind die Zahlen $k_1, k_2, \dots k_w$ sämmtlich gleich $n$, daher $D_k = 0$ und die Formel (7) geht über in $A_k + A_{n-k} = 2(\pi - 1) + w$. Die letztere Gleichung habe ich auf ganz anderem Wege in der Arbeit „über diejenigen algebraischen Gebilde, welche eindeutige Transformationen in sich zulassen“ (Nr.~6) erhalten.

\origpage{442}
Nun ist
\[
\begin{aligned}
A_0 + A_1 + \dots + A_{n-1} &= A_0 + \frac{1}{2} \sum_1^{n-1} {}_k (A_k + A_{n-k}) \\
&= \pi + \frac{1}{2} \sum_1^{n-1} {}_k [(2\pi - 2) + w - D_k] \\
&= \pi + (n-1)(\pi-1) + \frac{1}{2}(n-1)w - \frac{1}{2} \sum_1^{n-1} {}_k D_k.
\end{aligned} \tag{8}
\]
Gemäss der Bedeutung von $D_k$ finden wir den Werth von $\displaystyle\sum D_k$, wenn wir der Reihe nach die Anzahl der Glieder in der Zahlenreihe
\[
1, \ 2, \ 3, \ \dots \ n - 1 \tag{9}
\]
bestimmen, welche durch $k_1$, bez. $k_2$ \dots bez. $k_w$ theilbar sind, und alle diese Anzahlen addiren. In der Reihe (9) giebt es aber $\frac{n}{k_i} - 1$ Zahlen, welche durch $k_i$ theilbar sind, nämlich die Zahlen
\[
k_i, \ 2k_i, \ \dots \ \left(\frac{n}{k_i} - 1\right) k_i.
\]
Folglich ist
\[
\sum_1^{n-1} {}_k D_k = \frac{n}{k_1} + \frac{n}{k_2} + \dots + \frac{n}{k_w} - w.
\]
Tragen wir diesen Werth in (8) ein, so ergiebt sich nach leichter Umformung
\[
\begin{aligned}
A_0 + A_1 + \dots + A_{n-1} \\
= 1 + n(\pi - 1) + \frac{1}{2} \left[ n\left(1 - \frac{1}{k_1}\right) + n\left(1 - \frac{1}{k_2}\right) + \dots + n\left(1 - \frac{1}{k_w}\right) \right],
\end{aligned} \tag{10}
\]
und dieser Ausdruck erweist sich nach Formel (1) in Nr.~7 in der That gleich dem Geschlecht $p$ der Fläche $F$.

Königsberg i.~Pr., 10.~Februar 1892.
\end{document}
